%%%%%%%%%%%%%%%%%%%%%%%%%%%%%%%%%%%%%%%%%%%%%%%%%%%%%%%%%%%%%%%%%%%%%%%%%%%%%%%
% Harvard Academic Notes - 통합 마스터 템플릿
% 모든 강의 노트에 적용되는 통일된 스타일
% 버전: 2.1 - 가독성 개선 (선택적 최적화)
% 최종 수정일: 2025-11-17
%%%%%%%%%%%%%%%%%%%%%%%%%%%%%%%%%%%%%%%%%%%%%%%%%%%%%%%%%%%%%%%%%%%%%%%%%%%%%%%

\documentclass[11pt,a4paper]{article}

%========================================================================================
% 기본 패키지
%========================================================================================

% --- 한국어 지원 ---
\usepackage{kotex}

% --- 페이지 레이아웃 ---
\usepackage[top=20mm, bottom=20mm, left=20mm, right=18mm]{geometry}
\usepackage{setspace}
\onehalfspacing                      % 1.5배 줄간격
\setlength{\parskip}{0.5em}          % 문단 간격
\setlength{\parindent}{0pt}          % 들여쓰기 없음

% --- 표 관련 ---
\usepackage{booktabs}              % 고품질 표
\usepackage{tabularx}              % 자동 너비 조절 표
\usepackage{array}                 % 표 컬럼 확장
\usepackage{longtable}             % 여러 페이지 표
\renewcommand{\arraystretch}{1.1}  % 표 행간 조절

%========================================================================================
% 헤더 및 푸터
%========================================================================================

\usepackage{fancyhdr}
\pagestyle{fancy}
\fancyhf{}
\fancyhead[L]{\small\textit{BADM-572: Data-Driven Decision Making}}
\fancyhead[R]{\small\textit{Module 06}}
\fancyfoot[C]{\thepage}
\renewcommand{\headrulewidth}{0.5pt}
\renewcommand{\footrulewidth}{0.3pt}

% 첫 페이지는 헤더 없음
\fancypagestyle{firstpage}{
    \fancyhf{}
    \fancyfoot[C]{\thepage}
    \renewcommand{\headrulewidth}{0pt}
}

%========================================================================================
% 색상 정의 (파스텔 톤 + 다크모드 호환)
%========================================================================================

\usepackage[dvipsnames]{xcolor}

% 밝은 배경용 파스텔 색상
\definecolor{lightblue}{RGB}{220, 235, 255}      % 부드러운 파랑
\definecolor{lightgreen}{RGB}{220, 255, 235}     % 부드러운 초록
\definecolor{lightyellow}{RGB}{255, 250, 220}    % 부드러운 노랑
\definecolor{lightpurple}{RGB}{240, 230, 255}    % 부드러운 보라
\definecolor{lightgray}{gray}{0.95}              % 밝은 회색
\definecolor{lightpink}{RGB}{255, 235, 245}      % 부드러운 핑크
\definecolor{boxgray}{gray}{0.95}
\definecolor{boxblue}{rgb}{0.9, 0.95, 1.0}
\definecolor{boxred}{rgb}{1.0, 0.95, 0.95}

% 진한 색상 (테두리/제목용)
\definecolor{darkblue}{RGB}{50, 80, 150}
\definecolor{darkgreen}{RGB}{40, 120, 70}
\definecolor{darkorange}{RGB}{200, 100, 30}
\definecolor{darkpurple}{RGB}{100, 60, 150}

%========================================================================================
% 박스 환경 (tcolorbox) - 6가지 타입
%========================================================================================

\usepackage[most]{tcolorbox}
\tcbuselibrary{skins, breakable}

% 1. 개요 박스 (강의 시작 부분)
\newtcolorbox{overviewbox}[1][]{
    enhanced,
    colback=lightpurple,
    colframe=darkpurple,
    fonttitle=\bfseries\large,
    title=📚 강의 개요,
    arc=3mm,
    boxrule=1pt,
    left=8pt,
    right=8pt,
    top=8pt,
    bottom=8pt,
    breakable,
    #1
}

% 2. 요약 박스
\newtcolorbox{summarybox}[1][]{
    enhanced,
    colback=lightblue,
    colframe=darkblue,
    fonttitle=\bfseries,
    title=📝 핵심 요약,
    arc=2mm,
    boxrule=0.7pt,
    left=6pt,
    right=6pt,
    top=6pt,
    bottom=6pt,
    breakable,
    #1
}

% 3. 핵심 정보 박스
\newtcolorbox{infobox}[1][]{
    enhanced,
    colback=lightgreen,
    colframe=darkgreen,
    fonttitle=\bfseries,
    title=💡 핵심 정보,
    arc=2mm,
    boxrule=0.7pt,
    left=6pt,
    right=6pt,
    top=6pt,
    bottom=6pt,
    breakable,
    #1
}

% 4. 주의사항 박스
\newtcolorbox{warningbox}[1][]{
    enhanced,
    colback=lightyellow,
    colframe=darkorange,
    fonttitle=\bfseries,
    title=⚠️ 주의사항,
    arc=2mm,
    boxrule=0.7pt,
    left=6pt,
    right=6pt,
    top=6pt,
    bottom=6pt,
    breakable,
    #1
}

% 5. 예제 박스
\newtcolorbox{examplebox}[1][]{
    enhanced,
    colback=lightgray,
    colframe=black!60,
    fonttitle=\bfseries,
    title=📖 예제,
    arc=2mm,
    boxrule=0.7pt,
    left=6pt,
    right=6pt,
    top=6pt,
    bottom=6pt,
    breakable,
    #1
}

% 6. 정의 박스
\newtcolorbox{definitionbox}[1][]{
    enhanced,
    colback=lightpink,
    colframe=purple!70!black,
    fonttitle=\bfseries,
    title=📌 정의,
    arc=2mm,
    boxrule=0.7pt,
    left=6pt,
    right=6pt,
    top=6pt,
    bottom=6pt,
    breakable,
    #1
}

% 7. 중요 박스 (importantbox - warningbox와 유사)
\newtcolorbox{importantbox}[1][]{
    enhanced,
    colback=boxred,
    colframe=red!70!black,
    fonttitle=\bfseries,
    title=⚠️ 매우 중요,
    arc=2mm,
    boxrule=0.7pt,
    left=6pt,
    right=6pt,
    top=6pt,
    bottom=6pt,
    breakable,
    #1
}

% 8. cautionbox (warningbox와 동일)
\let\cautionbox\warningbox
\let\endcautionbox\endwarningbox

%========================================================================================
% 코드 블록 설정 (밝은 배경)
%========================================================================================

\usepackage{listings}

\definecolor{codegray}{rgb}{0.5,0.5,0.5}
\definecolor{codepurple}{rgb}{0.58,0,0.82}
\definecolor{backcolour}{rgb}{0.95,0.95,0.95}

\lstset{
    basicstyle=\ttfamily\small,
    backgroundcolor=\color{lightgray},
    keywordstyle=\color{darkblue}\bfseries,
    commentstyle=\color{darkgreen}\itshape,
    stringstyle=\color{purple!80!black},
    numberstyle=\tiny\color{black!60},
    numbers=left,
    numbersep=8pt,
    breaklines=true,
    breakatwhitespace=false,
    frame=single,
    frameround=tttt,
    rulecolor=\color{black!30},
    captionpos=b,
    showstringspaces=false,
    tabsize=2,
    xleftmargin=15pt,
    xrightmargin=5pt,
    escapeinside={\%*}{*)}
}

% Python 코드 스타일
\lstdefinestyle{pythonstyle}{
    language=Python,
    morekeywords={self, True, False, None},
}

% SQL 코드 스타일
\lstdefinestyle{sqlstyle}{
    language=SQL,
    morekeywords={SELECT, FROM, WHERE, JOIN, GROUP, BY, ORDER, HAVING},
}

%========================================================================================
% 목차 스타일링
%========================================================================================

\usepackage{tocloft}
\renewcommand{\cftsecleader}{\cftdotfill{\cftdotsep}}
\setlength{\cftbeforesecskip}{0.4em}
\renewcommand{\cftsecfont}{\bfseries}
\renewcommand{\cftsubsecfont}{\normalfont}

%========================================================================================
% 표 및 그림
%========================================================================================

\usepackage{graphicx}              % 이미지
\usepackage{adjustbox}             % 표/박스 크기 조절

% 표 캡션 스타일
\usepackage{caption}
\captionsetup[table]{
    labelfont=bf,
    textfont=it,
    skip=5pt
}
\captionsetup[figure]{
    labelfont=bf,
    textfont=it,
    skip=5pt
}

%========================================================================================
% 수학
%========================================================================================

\usepackage{amsmath, amssymb, amsthm}

% 정리 환경
\theoremstyle{definition}
\newtheorem{theorem}{정리}[section]
\newtheorem{lemma}[theorem]{보조정리}
\newtheorem{proposition}[theorem]{명제}
\newtheorem{corollary}[theorem]{따름정리}
\newtheorem{definition}{정의}[section]
\newtheorem{example}{예제}[section]

%========================================================================================
% 하이퍼링크
%========================================================================================

\usepackage[
    colorlinks=true,
    linkcolor=blue!80!black,
    urlcolor=blue!80!black,
    citecolor=green!60!black,
    bookmarks=true,
    bookmarksnumbered=true,
    pdfborder={0 0 0}
]{hyperref}

% PDF 메타데이터는 각 문서에서 설정
\hypersetup{
    pdftitle={BADM-572: Data-Driven Decision Making - Module 06},
    pdfauthor={강의 노트},
    pdfsubject={Academic Notes}
}

%========================================================================================
% 기타 유용한 패키지
%========================================================================================

\usepackage{enumitem}              % 리스트 커스터마이징
\setlist{nosep, leftmargin=*, itemsep=0.3em}

\usepackage{microtype}             % 타이포그래피 개선
\usepackage{footnote}              % 각주 개선
\usepackage{url}                   % URL 줄바꿈
\urlstyle{same}

%========================================================================================
% 사용자 정의 명령어
%========================================================================================

% 강조 텍스트
\newcommand{\important}[1]{\textbf{\textcolor{red!70!black}{#1}}}
\newcommand{\keyword}[1]{\textbf{#1}}
\newcommand{\term}[1]{\textit{#1}}
\newcommand{\code}[1]{\texttt{#1}}

% 용어 설명 (인라인)
\newcommand{\defterm}[2]{\textbf{#1}\footnote{#2}}

% 섹션 시작 전 페이지 분리
\newcommand{\newsection}[1]{\newpage\section{#1}}

%========================================================================================
% 문서 제목 스타일
%========================================================================================

\usepackage{titling}
\pretitle{\begin{center}\LARGE\bfseries}
\posttitle{\par\end{center}\vskip 0.5em}
\preauthor{\begin{center}\large}
\postauthor{\end{center}}
\predate{\begin{center}\large}
\postdate{\par\end{center}}

%========================================================================================
% 섹션 제목 간격
%========================================================================================

\usepackage{titlesec}
\titlespacing*{\section}{0pt}{1.5em}{0.8em}
\titlespacing*{\subsection}{0pt}{1.2em}{0.6em}
\titlespacing*{\subsubsection}{0pt}{1em}{0.5em}

%========================================================================================
% 메타 정보 박스 명령어
%========================================================================================

\newcommand{\metainfo}[4]{
\begin{tcolorbox}[
    colback=lightpurple,
    colframe=darkpurple,
    boxrule=1pt,
    arc=2mm,
    left=10pt,
    right=10pt,
    top=8pt,
    bottom=8pt
]
\begin{tabular}{@{}rl@{}}
▣ \textbf{강의명:} & #1 \\[0.3em]
▣ \textbf{주차:} & #2 \\[0.3em]
▣ \textbf{교수명:} & #3 \\[0.3em]
▣ \textbf{목적:} & \begin{minipage}[t]{0.75\textwidth}#4\end{minipage}
\end{tabular}
\end{tcolorbox}
}

%========================================================================================
% 끝
%========================================================================================


\begin{document}

\thispagestyle{firstpage}

\metainfo{BADM-572 데이터 기반 의사결정}{Module 2}{Fataneh Taghaboni-Dutta}{두 그룹 간 평균 차이의 통계적 유의미성 판단, 가설 검정 및 신뢰 구간 계산 방법을 학습하고, 비즈니스 맥락에서 두 가지 전략, 제품 또는 그룹의 효과를 비교하는 양측 및 단측 검정 방법을 습득합니다.}

\tableofcontents

\newpage

\begin{summarybox}[title=모듈 2: 추론 및 예측 통계 개요]
이 문서는 BADM-572 '데이터 기반 의사결정' 과목의 '추론 및 예측 통계 모듈 2'의 첫 번째 파트를 다룹니다.
주요 목표는 두 그룹 간의 평균 차이가 통계적으로 유의미한지 판단하는 가설 검정 절차와 신뢰 구간 계산 방법을 학습하는 것입니다.
특히, 비즈니스 맥락에서 두 가지 전략, 제품 또는 그룹의 효과를 비교하는 데 필요한 양측 및 단측 검정 방법을 다룹니다.
Excel을 활용한 실습을 통해 이론적 지식을 실제 데이터 분석에 적용하는 능력을 기르는 데 중점을 둡니다.
\end{summarybox}

\newpage
\section{용어 정리}
\addcontentsline{toc}{section}{용어 정리}

\begin{adjustbox}{width=\textwidth,center}
\begin{tabular}{lllc}
\toprule
\textbf{용어 (Term)} & \textbf{쉬운 설명} & \textbf{원어 (Original)} & \textbf{비고} \\
\midrule
가설 검정 & 표본 데이터를 사용하여 모집단에 대한 주장이 사실인지 통계적으로 확인하는 과정 & Hypothesis Testing & \\
귀무 가설 & 통계적 검정의 기본 가정으로, 보통 '차이가 없다'는 주장 & Null Hypothesis ($H_0$) & 기각 대상 \\
대립 가설 & 귀무 가설과 반대되는 주장으로, '차이가 있다'는 주장 & Alternative Hypothesis ($H_A$) & 채택 대상 \\
유의 수준 ($\alpha$) & 귀무 가설이 사실인데도 이를 기각할 확률의 최대 허용치 & Significance Level ($\alpha$) & 보통 0.05 (5\%) 또는 0.01 (1\%) \\
p-값 (p-value) & 귀무 가설이 사실일 때, 현재 관측된 표본 데이터 또는 그보다 극단적인 결과가 나타날 확률 & p-value & p-value < $\alpha$ 이면 귀무 가설 기각 \\
검정 통계량 (t-값) & 표본 데이터가 귀무 가설과 얼마나 다른지를 나타내는 표준화된 값 & Test Statistic (t-value) & \\
신뢰 구간 & 모집단 모수가 특정 범위 내에 있을 것이라고 추정되는 구간 & Confidence Interval & \\
단측 검정 & 한 방향으로의 차이(크거나 작다)를 검정하는 가설 검정 & One-Tail Test & \\
양측 검정 & 양방향으로의 차이(다르다)를 검정하는 가설 검정 & Two-Tail Test & \\
표준 오차 & 표본 통계량(예: 표본 평균)이 모집단 모수와 얼마나 차이 나는지 나타내는 척도 & Standard Error & \\
자유도 (df) & 통계량을 계산할 때 자유롭게 변할 수 있는 독립적인 값의 수 & Degrees of Freedom (df) & \\
\bottomrule
\end{tabular}
\end{adjustbox}

\newpage
\section{1. 서문 (Preface)}
\addcontentsline{toc}{section}{1. 서문 (Preface)}

Gies eBook을 선택해 주셔서 감사합니다.

이 Gies eBook은 Coursera에서 Fataneh Taghaboni-Dutta 교수님의 '비즈니스를 위한 추론 및 예측 통계' 모듈 2의 확장된 비디오 강의 스크립트를 기반으로 합니다. Gies eBook은 MOOC 비디오의 모든 정보를 완전히 접근 가능한 형식으로 제공하는 독서 경험을 제공합니다. 이 Gies eBook은 ePUB 3.0 형식을 지원하는 모든 표준 기반 전자책 리더 소프트웨어와 함께 사용할 수 있습니다.

각 Gies eBook은 레슨별로 나뉘어 있으며, 전자책 리더의 목차 기능을 사용하여 탐색할 수 있습니다. 각 레슨 내에서는 항상 다음 콘텐츠 순서가 나타납니다:
\begin{itemize}
    \item \textbf{레슨 제목 (Lesson title)}
    \item 각 레슨의 웹 기반 비디오 링크 (온라인 상태여야 시청 가능)
\end{itemize}
레슨 내에서 슬라이드가 변경되거나 다음 정보 비디오 장면으로 전환될 때마다 다음이 제공됩니다:
\begin{itemize}
    \item 현재 슬라이드 또는 비디오 장면의 \textbf{썸네일 이미지 (Thumbnail image)}
    \item 비디오 슬라이드에 있는 모든 텍스트는 썸네일 아래에 검색 가능하고 스크린 리더가 읽을 수 있는 형식으로 재현됩니다.
    \item 슬라이드에 제시된 그래프 및 차트와 같은 중요한 시각 자료에 대한 \textbf{확장된 텍스트 설명 (Extended text description)}
    \item 비디오의 모든 표 형식 데이터는 스크린 리더 탐색 및 읽기를 위해 재현되고 적절하게 레이블링됩니다.
    \item 모든 수학 방정식은 화면에 표시될 경우 내용과 표현을 모두 제공하는 MathML로 제시됩니다.
    \item 말하는 사람별로 레이블링된 비디오의 모든 원본 음성을 캡처한 \textbf{스크립트 (Transcript)}
\end{itemize}
모든 Gies eBook은 접근성과 사용성을 최우선으로 설계되었습니다. 이 디자인은 독자들이 선택하는 디지털 독서 도구에 관계없이 유연한 방식으로 모든 독자에게 서비스를 제공하기 위함입니다.

이 Gies eBook에 대한 질문이나 개선 제안이 있으시면 Giesbooks@illinois.edu로 문의하십시오.

\begin{infobox}[title=저작권 정보]
Copyright \textcopyright 2019 by Fataneh Taghaboni-Dutta \\
All rights reserved. \\
Published by the Gies College of Business at the University of Illinois at Urbana-Champaign, and the Board of Trustees of the University of Illinois
\end{infobox}

\newpage
\section{2. 모듈 2: 비즈니스를 위한 추론 및 예측 통계}
\addcontentsline{toc}{section}{2. 모듈 2: 비즈니스를 위한 추론 및 예측 통계}

이 모듈에서는 비즈니스 의사결정을 위한 추론 및 예측 통계의 핵심 개념을 탐구합니다. 특히, 두 그룹 간의 평균 또는 비율 차이를 통계적으로 검정하는 방법에 초점을 맞춥니다.

\subsection{2.1. 레슨 2-1: 평균 차이 없음 검정 (Testing for No Difference in Means)}
\addcontentsline{toc}{subsection}{2.1. 레슨 2-1: 평균 차이 없음 검정}

이 레슨에서는 두 모집단의 평균이 서로 다른지 여부를 판단하는 가설 검정 방법을 배웁니다. 이는 비즈니스에서 두 가지 전략, 제품 또는 그룹의 효과를 비교할 때 매우 유용합니다.

\subsubsection*{2.1.1. 레슨 2-1.1: 평균 차이 없음 검정}
\addcontentsline{toc}{subsubsection}{2.1.1. 레슨 2-1.1: 평균 차이 없음 검정}

\begin{summarybox}[title=핵심 요약: 두 평균의 차이 검정]
두 그룹 간의 평균 차이가 통계적으로 유의미한지 판단하는 가설 검정 절차를 학습합니다. 이는 단순히 표본 평균이 다르다고 해서 모집단 평균이 다르다고 단정할 수 없기 때문에 중요합니다.
\end{summarybox}

\begin{examplebox}[title=생각해 볼 질문: 실제 비즈니스에서의 차이]
\begin{itemize}
    \item 운동을 조금이라도 하면 전혀 운동하지 않는 것보다 더 오래 살 수 있을까요? (건강 분야)
    \item 고객에게 현금 환급 보상을 제공하는 것이 무료 여행 포인트를 제공하는 것보다 신용카드 사용을 더 늘릴까요? (마케팅/금융 분야)
\end{itemize}
\end{examplebox}

지난 모듈에서는 단일 평균 또는 비율을 고정된 값과 비교하는 방법을 살펴보았습니다. 이 모듈에서는 이를 확장하여 두 평균 또는 두 비율을 서로 비교합니다. 대부분의 상황에서 우리는 어떤 일을 다르게 하는 것이 아무것도 하지 않는 것과 비교하여 결과에 변화를 가져올지 궁금해합니다.

분명히 우리 중 일부, 아니 대부분은 건강상의 이점 때문에 운동을 합니다. 하지만 어떻게 그런 믿음을 갖게 되었을까요? 과학자들이 운동하는 사람들과 운동하지 않는 사람들의 건강과 수명을 비교했기 때문입니다. 또는 회사들이 한 가지 프로모션 전략을 다른 전략보다 사용하기로 결정하는 이유는 무엇일까요? 각 전략 하의 결과를 비교하여 한 전략이 다른 전략보다 우수하다고 생각할 만한 근거를 찾았기 때문입니다. 이것이 우리가 이 모듈에서 배울 내용입니다. 기본적으로 우리는 그룹 간의 비교를 올바르게 수행하는 방법과 그룹 간에 유의미한 차이가 있는지 여부를 결정하는 방법을 배울 것입니다.

\begin{figure}[h!]
    \centering
    % 이미지 파일이 제공되지 않았으므로 텍스트로 대체합니다.
    % % \includegraphics[width=0.8\textwidth]{example-morrow-plots.png}  % Image not found: example-morrow-plots.png
    \textbf{[이미지: UIUC의 모로우 플롯 (옥수수 밭 사진)]}
    \caption{UIUC의 모로우 플롯 (Morrow Plots at UIUC)}
    \label{fig:morrow_plots}
\end{figure}
\begin{infobox}[title=모로우 플롯: 1세기의 학습과 통계의 뿌리]
이곳 일리노이 대학교에서 우리가 배우려는 내용이 깊은 뿌리를 가지고 있다는 것을 자랑스럽게 말씀드립니다. 제 사무실 창문에서 모로우 플롯을 볼 수 있는데, 여름철에는 자라는 옥수수 밭이 이 그림에서 보시는 것처럼 아름답습니다.

1876년, 조지 모로우(George Morrow)와 또 다른 연구원 맨리 마일스(Manley Miles)는 이 플롯을 사용하여 실험을 수행했습니다. 그들은 옥수수 수확량을 늘리는 방법을 알고 싶었습니다. 그들은 일부 지역에는 비료를 주고 다른 지역에는 주지 않은 다음, 둘을 비교하여 더 나은 수확을 위해 어떤 종류의 비료를 얼마나 사용해야 하는지 알아냈습니다. 그들은 다른 작물에 대해서도 실험을 반복했습니다. 그들은 더 많은 실험을 하고 모든 결과를 통계를 사용하여 분석했으며, 그들의 연구를 통해 농업을 크게 개선한 많은 중요한 발견이 나왔습니다.

그들이 한 모든 것은 식물 과학과 통계라는 두 가지 과학을 기반으로 했습니다. 1800년대에는 통계가 주로 농업 분야에서 사용되었으며, 그들이 1세기 이상 전에 사용했던 통계 분석은 오늘날 다른 분야에도 동일하게 적용됩니다. 그럼, 두 그룹 간의 비교를 어떻게 수행하는지 살펴보겠습니다.
\end{infobox}

\begin{infobox}[title=평가 방법: '실제 차이'의 중요성]
우리는 표본 평균의 관측된 차이가 \textbf{실제 차이 (real difference)}를 나타낸다는 것을 입증해야 합니다.

\textbf{왜 중요할까요?} \\
여러분은 이제 한 모집단에서 표본을 추출하고 동일한 모집단에서 다른 표본을 추출하면 두 표본의 평균이 다를 가능성이 높다는 것을 알고 있습니다. 그렇다면 두 개의 다른 모집단에서 표본을 추출할 때 그들 사이에 차이가 나타나는 것은 놀라운 일이 아닙니다. 여기서 진짜 질문은 '언제 그 차이가 실제 차이이며, 단순히 무작위적이고 허용 가능한 자연적 변동이 아닌가?'입니다. 우리는 가설 검정을 수행하고 주어진 유의 수준에서 관측된 차이를 평가함으로써 이를 수행합니다.
\end{infobox}

\begin{infobox}[title=두 평균에 대한 가설 검정 절차]
두 모집단의 평균을 표본을 통해 비교할 때, 가설 검정 과정은 지난 모듈에서 보았던 것과 동일합니다. 즉, 다음 단계를 따릅니다:
\begin{enumerate}
    \item \textbf{$H_0$ (귀무 가설)와 $H_A$ (대립 가설)를 설정합니다.}
    \item \textbf{유의 수준 $\alpha$를 지정합니다.}
    \item \textbf{표본 데이터에 대한 p-값을 계산합니다.}
    \item \textbf{$H_0$를 기각하거나 기각하지 않습니다.}
\end{enumerate}
\end{infobox}

\begin{examplebox}[title=예시 1: 배터리 수명 비교 - 문제 정의]
한 상점 브랜드 배터리가 더 비싼 전국 브랜드 배터리만큼 오래 지속된다고 주장합니다. 소비자 감시 단체는 상점 브랜드의 주장이 거짓이 아닌지 5\% 유의 수준에서 확인하고 싶어 합니다. 배터리는 정상적인 사용을 모방하는 장비로 계속 테스트되었고, 배터리가 더 이상 작동하지 않을 때까지의 경과 시간이 기록되었습니다.

\textbf{배경 설명:} \\
소비자 감시 단체(예: Consumer Reports)는 많은 소비재를 테스트하고 그 결과를 발표하여 소비자들이 새로운 제품을 구매할 때 최고의 가치를 찾도록 돕습니다. 이 예시에서는 배터리 수명에 대한 주장을 검증하는 시나리오를 다룹니다.
\end{examplebox}

\begin{infobox}[title=두 브랜드 비교 (1/8): 문제 정의 및 변수 설정]
배터리 두 브랜드의 평균 수명 차이

\textbf{명칭 설정:}
\begin{itemize}
    \item 상점 브랜드: A
    \item 전국 브랜드: B
    \item $\mu_A$: 브랜드 A의 평균 수명 (시간 단위)
    \item $\mu_B$: 브랜드 B의 평균 수명 (시간 단위)
\end{itemize}
여기서는 배터리가 지속되는 평균 시간에 관심이 있으므로, 문제를 두 모집단 간의 차이로 구성할 수 있습니다.
\end{infobox}

\begin{infobox}[title=두 브랜드 비교 (2/8): 1단계 - 가설 설정]
\textbf{주장 (Claim):} 상점 브랜드 배터리는 전국 브랜드 배터리만큼 오래 지속된다.
$\mu_A = \mu_B$

\textbf{가설:}
\begin{itemize}
    \item $H_0 : \mu_A - \mu_B = 0$ (귀무 가설: 두 브랜드의 평균 수명 차이가 없다)
    \item $H_A : \mu_A - \mu_B \neq 0$ (대립 가설: 두 브랜드의 평균 수명 차이가 있다)
\end{itemize}
상점 브랜드의 주장은 전국 브랜드만큼 오래 지속된다는 것입니다. 즉, $\mu_A$는 $\mu_B$와 같습니다. 이를 귀무 가설과 대립 가설로 다시 작성하면, 귀무 가설은 평균 차이가 0이고, 대립 가설은 차이가 0이 아니라는 것입니다. 단일 표본 가설 검정과 마찬가지로, 이것은 \textbf{양측 검정 (two-tail test)}입니다. 나중에 단측 검정(one-tail test)의 예시도 다룰 것입니다. 지금은 이 질문을 해결하기 위해 수집된 데이터를 살펴봐야 합니다.
\end{infobox}

\begin{infobox}[title=두 브랜드 비교 (3/8): 2단계 - 유의 수준 설정]
\textbf{가설:}
\begin{itemize}
    \item $H_0 : \mu_A - \mu_B = 0$
    \item $H_A : \mu_A - \mu_B \neq 0$
\end{itemize}
\textbf{유의 수준:} $\alpha = 0.05$ (감시 단체는 5\% 유의 수준에서 이를 확인하고자 합니다.)
\end{infobox}

\begin{infobox}[title=두 브랜드 비교 (4/8): 데이터 수집의 중요성]
\textbf{비교를 위한 데이터:} 두 모집단에서 무작위로 독립적으로 추출된 데이터

\textbf{데이터 수집의 중요성:} \\
3단계로 넘어가기 전에, 이 두 브랜드를 비교하기 위한 데이터가 필요합니다. 데이터 수집은 결코 사소한 작업이 아니라는 점을 상기시켜 드립니다. 두 개의 다른 브랜드를 비교할 때는 '사과와 사과를 비교'하는 것처럼 신중해야 합니다. 예를 들어, 테스트에 사용된 방법은 가능한 한 유사해야 하며, 각 그룹에서 테스트된 배터리 종류(예: AAA 또는 D)도 거의 동일해야 합니다.

각 표본의 관측치 수는 동일할 필요가 없습니다. 다른 크기의 표본을 가질 수 있습니다. 여기서는 기관이 테스트를 수행하므로 두 브랜드에 대해 동일한 표본 크기를 사용할 것입니다. 그러나 설문조사를 하는 경우, 동일한 수의 응답을 받지 못할 수도 있으며, 이는 괜찮습니다.

지금은 데이터가 두 모집단에서 무작위로 독립적으로 선택된 배터리를 기반으로 생성되었다고 가정합니다.
\end{infobox}

\begin{infobox}[title=두 브랜드 비교 (5/8): 3단계 - p-값 계산 (기본 원리)]
p-값은 추정치와 귀무 가설을 분리하는 표준 오차의 수(검정 통계량)를 알면 찾을 수 있습니다.

\textbf{단일 모집단의 경우 검정 통계량:}
$t = \frac{\bar{x} - \mu_0}{s/\sqrt{n}}$

이전과 마찬가지로, 단일 표본에 대한 가설 검정을 수행했을 때처럼, p-값은 우리가 찾은 표본 결과와 같은 결과를 찾을 확률이 될 것입니다. 이 확률(p-값)은 우리의 추정치와 가설화된 차이(검정 통계량, $t$로 표시)를 분리하는 표준 오차의 수를 알면 찾을 수 있습니다.
단일 모집단의 평균에 대한 가설 검정을 수행했을 때 t-값을 어떻게 계산했는지 기억하시나요? 두 모집단을 비교할 때도 동일한 구조를 사용합니다.
\end{infobox}

\begin{infobox}[title=두 브랜드 비교 (6/8): 3단계 - p-값 계산 (두 모집단)]
\textbf{단일 모집단의 경우 검정 통계량:}
$t = \frac{\bar{x} - \mu_0}{s/\sqrt{n}}$

\textbf{두 모집단을 비교하는 경우 검정 통계량:}
$t = \frac{(\bar{x}_1 - \bar{x}_2) - D_0}{\sqrt{\frac{s_1^2}{n_1} + \frac{s_2^2}{n_2}}}$

단일 평균 대신 두 표본 평균의 차이를 사용합니다. 잠시 시간을 내어 두 번째 방정식을 분석하고 첫 번째 방정식과 어떻게 관련되는지 살펴보겠습니다.
\begin{itemize}
    \item \textbf{분자 (Numerator):} 두 표본 평균의 차이($\bar{x}_1 - \bar{x}_2$)가 첫 번째 인수로 사용되며, 첫 번째 방정식의 단일 표본 평균을 대체합니다. 두 번째 인수는 $D_0$인데, 이는 두 모집단 간의 가설화된 차이이며, 첫 번째 방정식의 가설화된 평균을 대체합니다. 이 예시에서는 이 차이를 0으로 가설화했습니다.
    \item \textbf{분모 (Denominator):} 표준 오차의 측정값입니다. 표준 오차를 계산하는 것은 상당히 복잡할 수 있지만, 다행히 소프트웨어 프로그램은 몇 가지 명령으로 이를 수행할 수 있으므로 표준 오차 계산 방법에 집중하지 않겠습니다. 대신 Excel을 사용하여 답을 얻을 것입니다. 따라서 예시를 해결하기 위해 슬라이드의 Excel 출력을 사용할 것입니다. 물론 나중에 Excel 데모 비디오를 시청하여 여기에 표시된 출력을 얻는 방법을 확인할 수 있습니다.
\end{itemize}
두 모집단을 비교하는 논리적 단계는 하나의 모집단과 하나의 표본만 있었을 때와 거의 동일하다는 것을 알아주셨으면 합니다. 수학이 조금 더 복잡해질 뿐입니다.

이제 데이터를 사용하여 예시를 진행해 보겠습니다. Excel을 사용하여 이 두 브랜드의 데이터를 분석했습니다. 다음은 그 출력입니다. t-값은 1.547로 강조 표시되어 있습니다. Excel은 단측 및 양측 검정의 모든 값을 제공합니다. 주어진 문제에 적합한 값에 집중해야 합니다. 여기서는 양측 검정을 수행하므로 출력에서 해당 값에만 집중할 것입니다.

양측 검정의 p-값은 0.123입니다. 이 숫자들은 무엇을 의미할까요? 우리가 가진 것을 더 잘 이해해 봅시다.

\begin{table}[h!]
    \centering
    \caption{Excel 결과: t-검정 (두 표본, 분산 불균등 가정) - 브랜드 A/B}
    \label{tab:excel_output_battery}
    \begin{adjustbox}{width=\textwidth,center}
    \begin{tabular}{lcc}
        \toprule
        \textbf{} & \textbf{브랜드 A (Store Brand)} & \textbf{브랜드 B (National Brand)} \\
        \midrule
        평균 (Mean) & 9.992972973 & 9.927567568 \\
        분산 (Variance) & 0.053700353 & 0.276790247 \\
        관측치 (Observations) & 185 & 185 \\
        가설 평균 차이 (Hypothesized Mean Difference) & 0 & \\
        자유도 (df) & 253 & \\
        \textbf{t-통계량 (t-stat)} & \textbf{1.547461849} & \\
        P(T $\le$ t) 단측 (one-tail) & 0.061500928 & \\
        T 임계값 단측 (critical one-tail) & 1.650898678 & \\
        \textbf{P(T $\le$ t) 양측 (two-tail)} & \textbf{0.123001856} & \\
        T 임계값 양측 (critical two-tail) & 1.969384804 & \\
        \bottomrule
    \end{tabular}
    \end{adjustbox}
\end{table}
\end{infobox}

\begin{infobox}[title=두 브랜드 비교 (7/8): 결과 해석 및 주의점]
\textbf{Excel 결과:} (위 표와 동일한 내용)

이 결과에 따르면, 상점 브랜드인 브랜드 A의 평균 수명은 9.99시간이고, 전국 브랜드인 브랜드 B의 평균 수명은 9.92시간입니다. 와우, 더 저렴한 브랜드가 더 좋은 브랜드처럼 보입니다!

하지만 잠깐, 이것은 전국 브랜드와 변동성 때문에 나타난 것일 수도 있으니 아직 판단을 내리지 마세요. 관측된 차이가 가설화된 차이 0에서 얼마나 떨어져 있는지를 나타내는 t-값은 1.547입니다. 그렇다면 5\% 유의 수준에서 표본에서 관측된 것과 같은 결과를 관측할 확률은 얼마일까요? 양측 검정의 p-값은 0.123입니다. 이제 우리의 공식으로 돌아가 결정을 내려봅시다.
\end{infobox}

\begin{infobox}[title=두 브랜드 비교 (8/8): 4단계 - 귀무 가설 기각 여부 결정]
\textbf{가설:}
\begin{itemize}
    \item $H_0 : \mu_A - \mu_B = 0$
    \item $H_A : \mu_A - \mu_B \neq 0$
\end{itemize}
\textbf{유의 수준:} $\alpha = 0.05$
\textbf{p-값:} 0.123

\textbf{결정 규칙:}
p-값 ($0.123$) > $\alpha$ ($0.05$) $\rightarrow$ 귀무 가설을 기각하지 않습니다 (Don't Reject the Null Hypothesis).

\textbf{결론:} \\
p-값 0.123이 $\alpha$ 0.05보다 크므로, 우리는 귀무 가설을 기각하지 않을 것입니다.
이는 우리의 표본 데이터가 더 저렴한 상점 브랜드가 더 비싼 전국 브랜드보다 가치나 품질이 떨어진다고 생각하게 할 만한 결과를 내놓지 않았다는 것을 의미합니다. 이것은 소비자에게 좋은 소식입니다. 우리는 동일한 품질을 얻으면서 돈을 절약할 수 있습니다.

기억하시겠지만, 브랜드 A의 평균 수명은 브랜드 B보다 약간 높았습니다. 그러나 우리의 완전한 분석을 기반으로 할 때, 두 브랜드 사이에 의미 있는 차이는 없다고 판단합니다. 여기서 보이는 차이는 '노이즈'로 간주되며 통계적으로 유의미하지 않습니다. 즉, "우리 브랜드가 더 비싼 브랜드보다 더 좋았다"고 전국적으로 광고하지 마세요. 적어도 이 데이터 세트에서는 그렇지 않습니다.
\end{infobox}

\begin{examplebox}[title=연습 문제: 온라인 vs. 매장 판매 - 문제 상황]
한 상점의 관리자는 웹사이트를 통한 일일 평균 판매액(달러 기준)이 매장 판매액과 다른지 알고 싶어 합니다. 데이터가 수집되었고 5\% 유의 수준에서 분석이 수행되었습니다. 관리자가 결과를 이해하도록 도와주세요.

\textbf{시작:} 가설과 유의 수준을 설정하세요.
\end{examplebox}

\begin{infobox}[title=연습 문제: 가설은 무엇인가요? (선택)]
관리자는 웹사이트를 통한 일일 평균 판매액(달러 기준)이 매장 판매액과 다른지 알고 싶어 합니다. 데이터가 수집되었고 5\% 유의 수준에서 분석이 수행되었습니다。 관리자가 결과를 이해하도록 도와주세요.

\textbf{가설은 무엇인가요?}
\begin{enumerate}
    \item $H_0 : \mu_{online} - \mu_{inStore} \neq 0$ 그리고 $H_A = \mu_{online} - \mu_{inStore} = 0$
    \item \textbf{$H_0 : \mu_{online} - \mu_{inStore} = 0$ 그리고 $H_A = \mu_{online} - \mu_{inStore} \neq 0$}
    \item $H_0 : \mu_{online} - \mu_{inStore} \ge 0$ 그리고 $H_A = \mu_{online} - \mu_{inStore} < 0$
    \item $H_0 : \mu_{online} - \mu_{inStore} \le 0$
\end{enumerate}
\end{infobox}

\begin{infobox}[title=연습 문제: 풀이 (1/2) - 가설 및 유의 수준 설정]
관리자는 두 판매 채널이 다른지 여부에 관심이 있으므로, 귀무 가설과 대립 가설을 양측 검정으로 설정합니다.
\begin{itemize}
    \item $H_0 : \mu_{online} - \mu_{inStore} = 0$ (귀무 가설: 온라인 평균 판매액과 매장 평균 판매액은 차이가 없다)
    \item $H_A : \mu_{online} - \mu_{inStore} \neq 0$ (대립 가설: 온라인 평균 판매액과 매장 평균 판매액은 차이가 있다)
\end{itemize}
\textbf{유의 수준:} $\alpha = 0.05$
\end{infobox}

\begin{infobox}[title=연습 문제: 풀이 (2/2) - 결과 분석 및 결정]
다음은 관리자의 데이터 분석 결과입니다. 52주간의 주간 판매 데이터가 있습니다. 귀무 가설을 기각하시겠습니까, 아니면 기각하지 않으시겠습니까? 기각하거나 기각하지 않는다는 것이 무엇을 의미할까요?

\begin{table}[h!]
    \centering
    \caption{Excel 결과: t-검정 (두 표본, 분산 불균등 가정) - 온라인/매장 판매}
    \label{tab:excel_output_sales}
    \begin{adjustbox}{width=\textwidth,center}
    \begin{tabular}{lcc}
        \toprule
        \textbf{} & \textbf{온라인 (Online)} & \textbf{매장 (In Store)} \\
        \midrule
        평균 (Mean) & 863.7380769 & 969.4619231 \\
        분산 (Variance) & 17374.14664 & 81505.48446 \\
        관측치 (Observations) & 52 & 52 \\
        가설 평균 차이 (Hypothesized Mean Difference) & 0 & \\
        자유도 (df) & 72 & \\
        t-통계량 (t-stat) & -2.424494522 & \\
        P(T $\le$ t) 단측 (one-tail) & 0.008920545 & \\
        T 임계값 단측 (critical one-tail) & 1.666293696 & \\
        \textbf{P(T $\le$ t) 양측 (two-tail)} & \textbf{0.01784109} & \\
        T 임계값 양측 (critical two-tail) & 1.993463567 & \\
        \bottomrule
    \end{tabular}
    \end{adjustbox}
\end{table}

\textbf{결정:}
p-값 ($0.0178$) < 0.05 ($\alpha$) $\rightarrow$ 귀무 가설을 기각합니다 (Reject $H_0$).

\textbf{해석:} \\
이것은 양측 검정이므로, 양측 검정 분석 결과에 집중해야 합니다. p-값을 직접 $\alpha$ (0.05)와 비교하여 결정을 내리겠습니다. 여기서 p-값은 0.0178이며, 이는 $\alpha$ 0.05보다 작습니다. 따라서 귀무 가설을 기각합니다.

이는 온라인 판매와 매장 판매의 주간 평균 판매액이 동일하지 않다는 것을 의미합니다.

이 예시에서 우리는 흥미로운 사실을 발견했습니다. 온라인과 매장 판매의 주간 평균 판매액 사이에 차이가 있습니다. 이제 자연스러운 질문은 '그 차이가 얼마나 되는가?'일 것입니다. 대부분의 사람들은 단순히 차이가 있다는 것을 아는 것만으로는 만족하지 않습니다. 그들은 차이의 크기도 알고 싶어 합니다. 분명히 이것은 중요합니다. 1달러 차이인가요, 아니면 수천 달러 차이인가요? 그럼 이 질문에 어떻게 답할 수 있는지 살펴보겠습니다.
\end{infobox}

\begin{infobox}[title=차이는 얼마나 될까? (How Much is the Difference?) - 평균 차이 계산]
\textbf{Excel 결과:} (위 표와 동일한 내용, 평균 행 강조)

\textbf{관측된 평균 차이 (Difference):} $969.46 - 863.74 = \$105.72$

\textbf{해석:} \\
분석 출력에서 매장 평균 판매액이 온라인보다 약간 높은 $105.72$달러임을 알 수 있습니다. 하지만 다른 표본을 추출하면 다른 결과가 나올 것이라는 것을 알고 있습니다. 따라서 차이를 추정하는 더 좋은 방법은 \textbf{신뢰 구간 (Confidence Interval)}을 사용하는 것입니다. 여기 있는 결과를 바탕으로 95\% 신뢰 구간을 개발할 수 있습니다.
\end{infobox}

\begin{infobox}[title=평균 차이에 대한 신뢰 구간 (Confidence Interval for the Mean Differences) - 공식]
\textbf{단일 표본의 경우 신뢰 구간:}
표본 평균 $\pm$ 오차 한계 (Margin of error)
$\left[\bar{x} \pm t_{\frac{\alpha}{2}} \frac{s}{\sqrt{n}}\right]$

\textbf{두 표본의 경우 신뢰 구간:}
$\left[(\bar{x}_1 - \bar{x}_2) \pm t_{\frac{\alpha}{2}} \sqrt{\frac{s_1^2}{n_1} + \frac{s_2^2}{n_2}}\right]$

다시 한번, 두 표본이 있을 때 해야 할 일을 단일 표본에 대한 신뢰 구간을 만들 때 배운 것과 연결시키고 싶습니다. 단일 표본이 있을 때 신뢰 구간은 표본 평균을 취하고 오차 한계를 더하고 빼서 계산되었습니다. 오차 한계는 위 방정식으로 표현되었습니다.

두 표본의 경우에도 아이디어는 동일합니다. 즉, 관측된 차이에 오차 한계를 더하고 뺄 것입니다. 그리고 방정식은 이제 위와 같이 보일 것입니다. 둘 사이의 유사성이 보이시나요? 우리는 본질적으로 동일한 작업을 수행하고 있습니다. 이제 예시로 돌아가 관리자를 위한 95번째 백분위수 신뢰 구간을 계산해 보겠습니다.
\end{infobox}

\begin{infobox}[title=차이는 얼마나 될까? (1/5) - 신뢰 구간 계산 시작]
\textbf{Excel 결과:} (온라인/매장 판매 t-검정 표)

\textbf{신뢰 구간 공식:}
$\left[(\bar{x}_1 - \bar{x}_2) \pm t_{\frac{\alpha}{2}} \sqrt{\frac{s_1^2}{n_1} + \frac{s_2^2}{n_2}}\right]$

\textbf{값 대입 (초기):}
$(969.46 - 863.74) \pm (1.993) \sqrt{\frac{81505.48}{52} + \frac{17374.14}{52}} = 105.72 \pm (1.993 \times 43.6)$

모든 표기법에 대한 값을 동일한 출력에서 찾고 수학적 계산을 수행하여 신뢰 구간을 얻을 수 있습니다. 표기법에 값을 대입할 때 오른쪽 방정식을 주의 깊게 살펴보세요.
\end{infobox}

\begin{infobox}[title=차이는 얼마나 될까? (2/5) - 평균 차이 대입]
\textbf{Excel 결과:} (온라인/매장 판매 t-검정 표)

\textbf{신뢰 구간 공식 (평균 차이 강조):}
$\left[\textbf{(\underline{969.46} - \underline{863.74})} \pm t_{\frac{\alpha}{2}} \sqrt{\frac{s_1^2}{n_1} + \frac{s_2^2}{n_2}}\right]$

\textbf{값 대입:}
$(969.46 - 863.74) \pm (1.993) \sqrt{\frac{81505.48}{52} + \frac{17374.14}{52}} = 105.72 \pm (1.993 \times 43.6)$

먼저 각 표본의 평균 지출 값을 대입하여 평균 차이를 얻습니다. 참고로 어떤 값을 먼저 쓰든 상관없습니다.
\end{infobox}

\begin{infobox}[title=차이는 얼마나 될까? (3/5) - t-임계값 대입]
\textbf{Excel 결과:} (온라인/매장 판매 t-검정 표)

\textbf{신뢰 구간 공식 (t-임계값 강조):}
$\left[(\bar{x}_1 - \bar{x}_2) \pm \textbf{\underline{$t_{\frac{\alpha}{2}}$}} \sqrt{\frac{s_1^2}{n_1} + \frac{s_2^2}{n_2}}\right]$

\textbf{값 대입:}
$(969.46 - 863.74) \pm (1.993) \sqrt{\frac{81505.48}{52} + \frac{17374.14}{52}} = 105.72 \pm (1.993 \times 43.6)$

다음으로 $t_{\frac{\alpha}{2}}$ 값을 찾아야 합니다. 이는 Excel 출력에서 '임계값 (critical value)'이라고도 불립니다. 우리의 경우 95\% 신뢰 구간을 원하며, 양측 검정의 t-값을 사용하면 1.993이 됩니다.
\end{infobox}

\begin{infobox}[title=차이는 얼마나 될까? (4/5) - 표준 오차 계산]
\textbf{Excel 결과:} (온라인/매장 판매 t-검정 표)

\textbf{신뢰 구간 공식 (표준 오차 강조):}
$\left[(\bar{x}_1 - \bar{x}_2) \pm t_{\frac{\alpha}{2}} \textbf{\underline{$\sqrt{\frac{s_1^2}{n_1} + \frac{s_2^2}{n_2}}$}}\right]$

\textbf{값 대입:}
$(969.46 - 863.74) \pm (1.993) \sqrt{\frac{81505.48}{52} + \frac{17374.14}{52}} = 105.72 \pm (1.993 \times 43.6)$

이제 표준 오차에 대한 값입니다. 여기 표기법은 각 표본의 분산을 각 표본 크기로 나눈 다음, 이들을 더하고 그 합의 제곱근을 취하는 것입니다. 이 경우, 매장 판매의 분산을 표본 크기로 나누고, 온라인 판매의 분산을 표본 크기로 나눈 다음, 이들을 더하고 제곱근을 취합니다. 이제 최종 결과를 얻기 위한 연산을 진행할 수 있습니다.
\end{infobox}

\begin{infobox}[title=차이는 얼마나 될까? (5/5) - 최종 신뢰 구간 및 해석]
\textbf{Excel 결과:} (온라인/매장 판매 t-검정 표)

\textbf{계산:}
$105.72 \pm (1.993 \times 43.6) = 105.72 \pm 86.93$

\textbf{95\% 신뢰 구간:}
$[\$18.80, \$192.65]$

\textbf{해석:} \\
모든 것이 올바르게 수행되었다면, 평균 차이는 $105.72$달러이고 오차 한계는 $86.93$달러가 됩니다. 이는 매장 판매와 온라인 판매의 주간 평균 판매액에 대한 95\% 신뢰 구간이 $18.80$달러에서 $192.65$달러 사이의 값이라는 것을 의미합니다. 이 구간 내의 어떤 값이라도 두 판매 채널 간의 실제 차이에 대한 가능성 있는 값입니다.

이제 관리자는 이 결과를 보고 "확실히 차이는 있지만, 실제 차이가 약 $19$달러에 불과하다면 개인적으로는 그렇게 중요하다고 생각하지 않을 것"이라고 말할 수 있습니다. 이것이 제가 여러분이 주목했으면 하는 부분입니다. 단순히 유의미한 차이를 발견했다고 해서 반드시 실질적인 것을 발견했다는 의미는 아닙니다. 통계를 이해하는 의사결정자는 이제 자신의 통찰력을 적용하여 무엇을 할지 스스로 결정할 수 있습니다.
\end{infobox}

\newpage
\subsubsection*{2.1.2. 레슨 2-1.2: Excel에서 평균 차이 없음 검정}
\addcontentsline{toc}{subsubsection}{2.1.2. 레슨 2-1.2: Excel에서 평균 차이 없음 검정}

\begin{summarybox}[title=핵심 요약: Excel을 활용한 t-검정]
Excel의 '데이터 분석' 도구를 사용하여 두 표본 평균의 차이에 대한 t-검정을 수행하는 실제 절차를 학습합니다. 특히, 분산이 같지 않다고 가정하는 경우의 검정 방법을 다룹니다.
\end{summarybox}

\begin{examplebox}[title=Excel 실습 예시: 배터리 수명 비교 (재확인)]
한 상점 브랜드 배터리가 더 비싼 전국 브랜드 배터리만큼 오래 지속된다고 주장합니다. 소비자 감시 단체는 상점 브랜드의 주장이 거짓이 아닌지 5\% 유의 수준에서 확인하고 싶어 합니다. 배터리는 정상적인 사용을 모방하는 장비로 계속 테스트되었고, 배터리가 더 이상 작동하지 않을 때까지의 경과 시간이 기록되었습니다.

\textbf{목표:} 두 표본의 평균을 비교하는 가설 검정을 사용하여 이를 테스트합니다. 한 표본은 전국 브랜드에서, 다른 하나는 상점 브랜드에서 가져옵니다.
\end{examplebox}

\begin{infobox}[title=Excel 워크시트 (2/9): 가설 설정 (Excel 준비)]
\textbf{브랜드 정의:}
\begin{itemize}
    \item A: 상점 브랜드 (Store Brand)
    \item B: 전국 브랜드 (National Brand)
\end{itemize}
\textbf{가설:}
\begin{itemize}
    \item $H_0 : \mu_A = \mu_B \rightarrow \mu_A - \mu_B = 0$ (귀무 가설: 두 브랜드의 평균 수명은 동일하다)
    \item $H_A : \mu_A \neq \mu_B \rightarrow \mu_A - \mu_B \neq 0$ (대립 가설: 두 브랜드의 평균 수명은 다르다)
\end{itemize}
가설을 설정하는 쉬운 방법은 상점 브랜드와 전국 브랜드의 수명이 같다고 생각하는 것입니다. 따라서 대립 가설은 다르다는 것이 됩니다. 이렇게 작성하면 소프트웨어가 처리할 수 있는 방식으로 작성하기가 더 쉽습니다.

소프트웨어는 두 값 사이의 차이가 무엇이라고 생각하는지 찾고 있습니다. 따라서 이제 이들을 '두 배터리 수명 간의 차이가 0과 같다'는 귀무 가설과 '차이가 0이 아니다'는 대립 가설로 다시 작성할 수 있습니다. 가설이 설정되면 단측 검정인지 양측 검정인지 알 수 있습니다. 이것은 양측 검정의 예시입니다. 왜냐하면 상점 브랜드가 전국 브랜드보다 나쁘거나 좋다고 말하는 것이 아니라, 단순히 '같다'고 말하기 때문입니다. 따라서 우리는 이 가설을 기각할 것입니다. 만약 상점 브랜드가 실제로 더 좋거나 더 나쁘다는 것을 발견한다면 말이죠. 즉, 우리가 예상하는 0이라는 차이에서 너무 많이 벗어나면 귀무 가설을 기각하게 됩니다. 그렇지 않으면 기각하지 않을 것입니다.
\end{infobox}

\begin{infobox}[title=Excel 워크시트 (3/9): 데이터 준비 및 표본 크기]
\textbf{데이터:} 배터리 수명 Excel 파일 (첫 번째 워크시트 참조)

\textbf{데이터 설명:} \\
데이터를 살펴보면 브랜드 A와 브랜드 B에 대해 각각 186개의 관측치가 있습니다. 즉, 동일한 수의 배터리가 테스트되었습니다. 참고로, 표본 크기가 동일할 필요는 없습니다. 이들은 두 개의 독립적인 모집단에서 무작위로 선택된 두 개의 독립적인 표본입니다. 따라서 동일할 필요는 없습니다. 배터리를 테스트하는 경우 동일한 수를 선택할 가능성이 높지만, 예를 들어 두 가지 다른 유형의 모집단에 설문조사를 보냈다면 한 모집단이 다른 모집단보다 더 많은 설문지를 반환할 수 있습니다. 따라서 표본 크기가 동일할 필요가 없다는 것은 우리에게 큰 장점입니다.
\end{infobox}

\begin{infobox}[title=Excel 워크시트 (4/9): Excel에서 t-검정 선택 절차]
\textbf{절차:}
\begin{enumerate}
    \item Excel에서 \textbf{데이터 (Data)} 탭으로 이동합니다.
    \item \textbf{데이터 분석 (Data Analysis)}을 선택합니다. (이 기능이 보이지 않으면 Excel 추가 기능에서 '분석 도구 팩'을 활성화해야 합니다.)
    \item 팝업 창이 나타나면 아래로 스크롤하여 세 가지 다른 t-검정을 찾습니다.
    \item \textbf{t-검정: 두 표본 평균 차이 검정 (분산 불균등 가정) (t-Test: Two-Sample Assuming Unequal Variances)}을 선택합니다.
\end{enumerate}

\textbf{분산 가정에 대한 설명:} \\
\begin{itemize}
    \item \textbf{분산 동일 가정 (Assuming Equal Variances):} 모집단 A(브랜드 A)와 모집단 B(브랜드 B) 내부에 변동성이 있지만, 그 변동성이 거의 같다고 가정하는 것입니다. 수동으로 계산할 때는 자유도를 구하기가 훨씬 간단해지기 때문에 이런 가정을 사용했습니다.
    \item \textbf{분산 불균등 가정 (Assuming Unequal Variances):} 이 가정이 실제와 다를 수 있습니다. 컴퓨터 프로그램을 사용할 때는 수동으로 어려운 계산을 할 필요가 없으므로, 분산이 같다고 가정할 이유가 거의 없습니다.
\end{itemize}
\textbf{권장 사항:} 항상 \textbf{분산 불균등 가정}으로 이 종류의 검정을 실행하세요. 분산이 실제로 동일하더라도 결과는 같을 것입니다. 하지만 분산이 동일하다고 가정했는데 실제로는 동일하지 않다면, 분석이 잘못될 수 있습니다. 따라서 안전하게 가장 일반화된 형태를 선택하세요. 소프트웨어를 사용할 때는 계산이 순식간에 이루어지므로 차이를 전혀 느끼지 못할 것입니다.

따라서 t-검정을 수행할 것입니다. 이 t-검정은 유의미한 차이가 존재하는지 여부를 알려줄 것입니다. 분산 불균등 가정을 사용하는 두 표본 t-검정을 선택하고 '확인'을 클릭합니다.
\end{infobox}

\begin{infobox}[title=Excel 워크시트 (5/9): t-검정 설정 입력]
\textbf{팝업 창 설정:}
\begin{itemize}
    \item \textbf{변수 1 범위 (Variable 1 Range):} \$A\$1:\$A\$186 (브랜드 A 데이터 선택)
    \item \textbf{변수 2 범위 (Variable 2 Range):} \$B\$1:\$B\$186 (브랜드 B 데이터 선택)
    \textbf{팁:} A1에 커서를 놓고 Ctrl+Shift+아래쪽 화살표를 눌러 전체 데이터 세트를 선택합니다.
    \item \textbf{레이블 (Labels):} 체크박스 선택 (변수 이름이 포함된 경우)
    \textbf{주의:} 이 옵션을 선택하지 않으면 "숫자가 아닌 값이 있습니다" 오류가 발생할 수 있습니다.
    \item \textbf{가설 평균 차이 (Hypothesized Mean Difference):} 0 (또는 비워둠)
    \item \textbf{알파 ($\alpha$):} 0.05 (기본값, 필요시 0.10 등으로 변경 가능)
    \item \textbf{출력 옵션 (Output Options):} 출력 범위 (Output range) 선택 후 원하는 셀 지정 (예: \$E\$3)
    \textbf{팁:} 출력 범위를 지정하기 전에 해당 옵션을 클릭해야 커서가 점프하여 드롭다운 메뉴가 회색으로 변하는 것을 방지할 수 있습니다.
\end{itemize}

\textbf{설명:} \\
"범위 1은 무엇입니까?"라고 묻습니다. 저는 레이블을 포함하는 것을 선호합니다. 그러면 분석에 레이블이 변수 1로 표시되기 때문입니다. A1에 커서를 놓고 Ctrl+Shift+아래쪽 화살표를 눌러 전체 데이터 세트를 선택합니다. 다시 위로 스크롤하여 브랜드 B에 커서를 놓고 전체 셀을 선택합니다. 레이블이 있다고 말하는 것을 잊지 마세요. 만약 잊으면 "숫자가 아닌 값이 있습니다"라는 오류 메시지가 나타날 것입니다. 그런 경우 다시 열어서 레이블을 클릭하면 됩니다.

"가설 평균 차이는 무엇입니까?"라고 묻습니다. 여기에 0을 입력하거나 무시해도 됩니다. $\alpha$는 0.05이며, 이것은 기본값입니다. 0.10 등으로 변경할 수 있습니다. 저는 출력을 같은 페이지에 두는 것을 선호하지만, 여기를 클릭하면 커서가 점프하여 이 드롭다운 메뉴가 회색으로 변합니다. 따라서 스프레드시트의 아무 곳이나 클릭하기 전에 여기를 클릭하고, 위치를 선택한 다음 '확인'을 클릭해야 합니다.
\end{infobox}

\begin{infobox}[title=Excel 워크시트 (6/9): 결과 출력 및 초기 해석]
\textbf{Excel 출력:} (이전 슬라이드의 배터리 t-검정 표와 동일)

\textbf{결과 서식 지정:} \\
출력이 뭉개져 보일 수 있으므로, 먼저 읽기 쉽게 서식을 지정합니다. 가장 좋은 방법은 커서를 여기에 놓고 두 번 클릭하면 크기가 늘어납니다. 걱정할 필요 없습니다.

\textbf{기술 통계량 (Descriptive Statistics):} \\
출력에서 가장 먼저 볼 수 있는 것은 많은 기술 통계량입니다. 예를 들어, 브랜드 A의 평균은 9.99시간이고 브랜드 B의 평균은 9.92시간이라고 알려줍니다. 브랜드 A는 상점 브랜드이고 브랜드 B는 전국 브랜드입니다. 기억하기 쉽게 바꾸고 싶다면 그렇게 할 수 있습니다. 예를 들어, '상점 브랜드'로 변경하면 생각하기 더 쉬울 것입니다.

이것을 보면 브랜드 A가 실제로 약간 더 나은 것처럼 보입니다. 이것은 상점 브랜드이므로 더 저렴한 배터리, 즉 일반 브랜드입니다. 그래서 "와, 브랜드 A가 실제로 더 좋네"라고 말할 수도 있습니다. 하지만 기억하세요, 우리는 그렇게 말할 수 없습니다. 이 경우 실제로는 같을 수도 있기 때문입니다. 그것이 바로 여기서 테스트하는 것입니다. 또한 브랜드 A의 분산과 브랜드 B의 분산도 알려줍니다. 분산은 표준 편차의 제곱이라는 것을 기억하세요. 따라서 표준 편차를 원한다면 이 값의 제곱근을 취하면 브랜드 A의 표준 편차를 알 수 있습니다.

브랜드 A에는 185개의 관측치가 있고 브랜드 B에도 185개의 관측치가 있다고 합니다. 가설 평균 차이는 0이었고, 자유도는 이 두 값을 사용하여 분산 불균등 가정에 사용되는 방정식을 기반으로 자유도를 생성하는 과정을 거쳐 계산되었습니다.

\textbf{t-통계량 및 p-값:} \\
이것이 완료되면 t-통계량 값과 단측 검정 및 양측 검정에 대한 p-값을 볼 수 있습니다. 또한 t-임계값이 어디에 있는지도 알려줍니다. t-임계값의 의미를 보여드리겠습니다.
\end{infobox}

\begin{infobox}[title=Excel 워크시트 (7/9): t-임계값 시각화 및 이해]
\textbf{t-임계값 (T critical value):}
\begin{itemize}
    \item 양측 검정 (Two-tail test)에서 유의 수준이 5\%인 경우, 왼쪽 꼬리에 2.5\%, 오른쪽 꼬리에 2.5\%가 있습니다.
    \item 이 경우 t-임계값은 1.96과 -1.96이 됩니다. (Excel 출력에서는 1.969384804)
\end{itemize}
\begin{figure}[h!]
    \centering
    % 이미지 파일이 제공되지 않았으므로 텍스트로 대체합니다.
    % % \includegraphics[width=0.7\textwidth]{example-bell-curve-two-tail.png}  % Image not found: example-bell-curve-two-tail.png
    \textbf{[이미지: 종 모양 곡선, 중앙에 0, 양쪽에 -1.96과 1.96으로 표시된 임계값, 각 꼬리 영역에 2.5\% 표시]}
    \caption{양측 검정의 t-임계값과 기각 영역}
    \label{fig:bell_curve_two_tail}
\end{figure}
\end{infobox}

\begin{infobox}[title=Excel 워크시트 (8/9): p-값 기반 최종 결정]
\textbf{Excel 출력:} (이전 슬라이드의 배터리 t-검정 표와 동일, P(T $\le$ t) 양측 값 0.1230 강조)

\textbf{결정:}
p-값 (양측) = 0.1230 > 0.05 ($\alpha$) $\rightarrow$ 귀무 가설을 기각하지 않습니다.

\textbf{설명:} \\
만약 단측 검정을 수행했다면, 어느 꼬리든 상관없이(오른쪽 꼬리든 왼쪽 꼬리든) 전체 5\%가 한쪽에 있을 것입니다. 예를 들어 오른쪽 꼬리 검정이라면 1.65가 될 것입니다 (Excel 출력에서는 1.650898678). Excel은 단측 검정인지 양측 검정인지 모르기 때문에 모든 값을 제공합니다. 여러분이 올바른 것을 선택해야 합니다.

또한 t-통계량도 알려줍니다. t-통계량(1.547)은 브랜드 A에서 표본을, 브랜드 B에서 표본을 추출했을 때, 가설화된 차이가 0일 때 이 두 표본에서 관찰된 차이가 어디쯤에 있는지를 말해줍니다. 그리고 p-값을 제공하며, 여러분은 가설을 기각할지 말지를 결정해야 합니다.

우리 사례에서는 양측 검정을 수행하고 있으므로, 테이블의 이 부분에만 집중할 것입니다. 이 p-값을 보면 0.05보다 크다는 것을 알 수 있습니다. 따라서 귀무 가설을 기각하지 않습니다. 귀무 가설을 기각하지 않는다는 것은 무엇을 의미할까요?
\end{infobox}

\begin{infobox}[title=Excel 워크시트 (9/9): 최종 결론 도출]
\textbf{가설:}
\begin{itemize}
    \item $H_0 : \mu_A - \mu_B = 0$ (강조)
    \item $H_A : \mu_A - \mu_B \neq 0$
\end{itemize}
\textbf{결론:} \\
우리의 가설은 이 두 브랜드가 거의 동일하게 작동한다는 것이었고, 데이터를 기반으로 이를 기각하지 않았습니다. 따라서 두 브랜드는 거의 동일합니다. 현재 그들이 주장하는 바는 우리가 가진 데이터를 기반으로 반박할 수 없습니다.
\end{infobox}

\newpage
\subsubsection*{2.1.3. 레슨 2-1.3: 두 표본에 대한 양측 평균 검정 및 Excel에서 신뢰 구간 계산}
\addcontentsline{toc}{subsubsection}{2.1.3. 레슨 2-1.3: 두 표본에 대한 양측 평균 검정 및 Excel에서 신뢰 구간 계산}

\begin{summarybox}[title=핵심 요약: 가설 검정 후 신뢰 구간 계산]
Excel의 데이터 분석 기능을 활용하여 두 표본 평균의 차이에 대한 가설 검정을 수행하고, 이어서 해당 차이에 대한 신뢰 구간을 계산하는 방법을 실습합니다. 이는 통계적 유의성뿐만 아니라 차이의 실제적인 크기를 이해하는 데 필수적입니다.
\end{summarybox}

\begin{examplebox}[title=Excel 실습 예시: 온라인 vs. 매장 판매 (신뢰 구간)]
한 상점의 관리자는 웹사이트를 통한 일일 평균 판매액(달러 기준)이 매장 판매액과 다른지 알고 싶어 합니다. 데이터가 수집되었고 5\% 유의 수준에서 분석이 수행되었습니다. 관리자가 결과를 이해하도록 도와주세요.

\textbf{목표:} 가설을 설정하고 유의 수준을 지정하여 분석을 시작합니다.
\end{examplebox}

\begin{infobox}[title=온라인/매장 Excel 워크시트 (2/10): 가설 설정]
\textbf{가설:}
\begin{itemize}
    \item $H_0 : \mu_{Online} = \mu_{Instore} \rightarrow \mu_{Online} - \mu_{Instore} = 0$ (귀무 가설: 온라인 판매와 매장 판매의 평균은 동일하다)
    \item $H_A : \mu_{Online} \neq \mu_{Instore} \rightarrow \mu_{Online} - \mu_{Instore} \neq 0$ (대립 가설: 온라인 판매와 매장 판매의 평균은 다르다)
\end{itemize}
관리자는 온라인과 매장 판매 사이에 차이가 있는지 궁금해합니다. 따라서 귀무 가설은 둘 사이에 차이가 없다는 것이고, 대립 가설은 차이가 있다는 것입니다. 이렇게 작성하면 우리가 무엇을 하고 있는지 더 쉽게 알 수 있습니다. 즉, 두 $\mu$의 차이가 0과 같거나 0과 같지 않다고 말하는 것입니다. 이제 데이터를 살펴보고 분석을 수행해 봅시다.
\end{infobox}

\begin{infobox}[title=온라인/매장 Excel 워크시트 (3/10): 데이터 준비]
\textbf{데이터:} 온라인 vs. 매장 Excel 파일 (첫 번째 워크시트 참조)

\textbf{데이터 설명:} \\
여기에 데이터가 있습니다. 이것은 여러 날 동안 온라인과 매장에서 기록된 일일 판매액(달러)입니다.
\end{infobox}

\begin{infobox}[title=온라인/매장 Excel 워크시트 (4/10): Excel에서 t-검정 선택]
\textbf{절차:}
\begin{enumerate}
    \item Excel에서 \textbf{데이터 (Data)} 탭으로 이동합니다.
    \item \textbf{데이터 분석 (Data Analysis)}을 선택합니다.
    \item \textbf{t-검정: 두 표본 평균 차이 검정 (분산 불균등 가정) (t-Test: Two-Sample Assuming Unequal Variances)}을 선택합니다.
    \item '확인'을 클릭합니다.
\end{enumerate}
다시 한번, 저는 분산 불균등 가정을 할 것입니다. 단순화하고 싶지 않습니다.
\end{infobox}

\begin{infobox}[title=온라인/매장 Excel 워크시트 (5/10): t-검정 설정 입력 (열 선택)]
\textbf{팝업 창 설정:}
\begin{itemize}
    \item \textbf{변수 1 범위 (Variable 1 Range):} \$A:\$A (온라인 판매 데이터 열 전체 선택)
    \item \textbf{변수 2 범위 (Variable 2 Range):} \$B:\$B (매장 판매 데이터 열 전체 선택)
    \textbf{팁:} Ctrl+Shift+아래쪽 화살표 대신 열 전체를 선택하는 것이 더 빠릅니다. 이 유형의 분석에서는 작동하지만 항상 이런 식으로 선택할 수 있는 것은 아닙니다.
    \item \textbf{레이블 (Labels):} 체크박스 선택
    \item \textbf{가설 평균 차이 (Hypothesized Mean Difference):} 0
    \item \textbf{알파 ($\alpha$):} 0.05
    \item \textbf{출력 옵션 (Output Options):} 새 워크시트 (New Worksheet Ply) 선택 (또는 출력 범위 지정)
\end{itemize}

\textbf{설명:} \\
"범위 1을 묻습니다. 이전에 Ctrl+Shift+아래쪽 화살표를 사용하여 선택하는 방법을 보여드렸지만, 여기서는 열 전체를 선택하는 것이 더 빠릅니다. 이 유형의 분석에서는 작동하지만 항상 이런 식으로 선택할 수 있는 것은 아닙니다. 하지만 여기서는 가능하며 선택 과정을 빠르게 할 수 있습니다.

두 번째 변수는 B열에 나타날 것입니다. 저는 차이가 없다고 가정합니다. 선입견이 없습니다. 레이블이 있다고 선택할 것입니다. 출력 범위를 선택하고, 여기를 클릭하고, 스프레드시트의 아무 곳이나 클릭한 다음 '확인'을 클릭합니다.

이제 결과가 나왔습니다. 다시 한번, 저는 양측 검정을 수행하고 있다는 것을 알고 있습니다. 왜냐하면 '같다' 대 '같지 않다'의 가설이기 때문입니다.
\end{infobox}

\begin{infobox}[title=온라인/매장 Excel 워크시트 (6/10): 결과 분석 및 결정]
\textbf{Excel 출력:} (온라인/매장 판매 t-검정 표)

\textbf{결정:}
p-값 (양측) = 0.017 < 0.05 ($\alpha$) $\rightarrow$ 귀무 가설을 기각합니다.

\textbf{설명:} \\
이 출력의 양측 부분에만 집중하면 됩니다. p-값을 보면 0.05보다 작으므로, 귀무 가설(두 판매 채널이 동일하다는 가설)을 기각하게 됩니다.

이제 두 판매 채널 사이에 차이가 있다는 것을 알았으니, 온라인 판매와 매장 판매 사이의 예상되는 차이가 무엇인지 알고 싶습니다. 이제 동일한 출력을 사용하여 두 판매 채널 사이에 존재하는 차이에 대한 신뢰 구간을 개발할 수 있습니다.
\end{infobox}

\begin{infobox}[title=온라인/매장 Excel 워크시트 (7/10): 신뢰 구간 계산 준비]
\textbf{Excel 출력:} (온라인/매장 판매 t-검정 표)

\textbf{신뢰 구간 공식:}
$\left[(\bar{x}_1 - \bar{x}_2) \pm t_{\frac{\alpha}{2}} \sqrt{\frac{s_1^2}{n_1} + \frac{s_2^2}{n_2}}\right]$

\textbf{평균 차이 (Mean Difference):}
$-105.7238462$ (온라인 평균 - 매장 평균)

\textbf{설명:} \\
두 판매 채널 간의 차이에 대한 신뢰 구간을 개발하기 위해, 이것은 두 채널 간의 평균 차이입니다. 하지만 이것만으로는 충분하지 않습니다. 단순히 이 둘의 차이를 취하여 이것이 예상되는 차이라고 말할 수는 없습니다. 충분히 큰 오차가 있습니다. 표본마다 변동성이 있습니다. 따라서 전체 차이를 개발해야 합니다.

지금 당장은 두 평균의 차이를 보고 이것이 $\bar{x}_1 - \bar{x}_2$라고 말할 수 있습니다. 이것은 이 값에서 이 값을 빼서 개발할 수 있습니다. 참고로, 어떤 것을 먼저 하든 상관없습니다. 부호는 바뀌겠지만, 차이의 절대값은 여전히 동일하며 그것이 우리가 신경 쓰는 부분입니다. 지금은 온라인 평균 판매액이 매장 판매액보다 105달러 적다고 말할 수 있습니다. 반대로 했다면, 이 양수는 매장 판매액이 온라인보다 105달러 많다고 말했을 것입니다. 따라서 실제로는 차이가 없습니다.

우리가 가진 차이에 대한 신뢰 구간 공식으로 돌아가서, 다음으로 찾아야 할 것은 t-값입니다. 신뢰 구간은 항상 양측입니다. 오차 한계는 플러스 또는 마이너스로 발생할 수 있으므로, 여기의 양측 t-값은 1.99입니다. 그것이 저의 t-값입니다. 다시 한번, 우리 공식을 보면 오차 한계를 구하기 위해 다음으로 필요한 것은 표준 오차입니다.
\end{infobox}

\begin{infobox}[title=온라인/매장 Excel 워크시트 (8/10): 표준 오차 계산]
\textbf{Excel 출력:} (온라인/매장 판매 t-검정 표)

\textbf{t-값:} $1.993463567$

\textbf{표준 오차 (Standard Error) 계산:}
\begin{lstlisting}[language=Excel, caption=Excel에서 표준 오차 계산]
=SQRT(E6/E7+F6/F7)
% E6: 온라인 분산 (17374.14)
% E7: 온라인 관측치 (52)
% F6: 매장 분산 (81505.48)
% F7: 매장 관측치 (52)
\end{lstlisting}
\textbf{결과:} $43.60655189$

\textbf{설명:} \\
우리 방정식에 기반한 표준 오차는 단순히 첫 번째 표본의 분산을 해당 표본 크기(이 경우 52)로 나눈 값에 두 번째 표본의 분산을 해당 표본 크기로 나눈 값을 더한 다음, 그 합의 제곱근을 취하는 것입니다. 두 값을 얻으면 제곱근을 사용하여 표준 오차를 찾을 수 있습니다. 이제 오차 한계를 계산할 준비가 되었습니다. 오차 한계는 단순히 t-값에 표준 오차를 곱한 것입니다.
\end{infobox}

\begin{infobox}[title=온라인/매장 Excel 워크시트 (9/10): 오차 한계 및 신뢰 구간 계산]
\textbf{Excel 출력:} (온라인/매장 판매 t-검정 표)

\textbf{계산된 값 요약:}
\begin{itemize}
    \item 평균 차이 (mean diff): $-105.7238462$
    \item t-값 (t value): $1.993463567$
    \item 표준 오차 (standard error): $43.60655189$
\end{itemize}
\textbf{오차 한계 (Margin of error):} $86.92807245$ ($1.993463567 \times 43.60655189$)

\textbf{신뢰 구간 계산:}
\begin{itemize}
    \item \textbf{하한 (Lower):} $-105.7238462 - 86.92807245 = -192.6519186$
    \item \textbf{상한 (Upper):} $-105.7238462 + 86.92807245 = -18.7957737$
\end{itemize}
\textbf{결론:} \\
우리는 95\% 신뢰 수준에서 온라인 판매와 매장 판매 간의 차이가 온라인 채널이 매장 판매보다 하루 평균 $18.80$달러에서 $192.65$달러 적게 판매할 수 있다는 것을 의미합니다.
\end{infobox}

\begin{infobox}[title=온라인/매장 Excel 워크시트 (10/10): 통계적 유의성과 실질적 의미의 차이]
\textbf{통계적 유의성 (Statistical Significance):}
\begin{itemize}
    \item 평균 차이 (mean diff): $-105.7238462$
    \item t-값 (t value): $1.993463567$
    \item 표준 오차 (standard error): $43.60655189$
    \item 오차 한계 (Margin of error): $86.92807245$
    \item 하한 (Lower): $-192.6519186$
    \item 상한 (Upper): $-18.7957737$
\end{itemize}

\textbf{해석:} \\
이것은 두 채널 간에 통계적으로 유의미한 차이를 보여줍니다. 하지만 실제로는 그 차이가 $18$달러만큼 작을 수도 있습니다. 따라서 관리자로서 여러분은 이것이 중요하지 않다고 판단할 수도 있습니다. 즉, 통계적으로 유의미하다고 해서 항상 의미 있는 것을 발견했다는 뜻은 아닙니다. 만약 이 숫자가 수천 달러의 차이였다면, 아무도 실제적으로도 중요한 문제라고 반박하지 않을 것입니다.

비즈니스 담당자로서 이것이 충분히 유의미하다고 판단하는 임계값은 해당 분야와 그 분야와 관련된 결정에 따라 달라지므로, 표준적인 답은 없습니다. 이것이 바로 여러분의 전문 지식이 발휘되는 지점이며, 단순히 통계적 값만 보는 것이 아닙니다.
\end{infobox}

\newpage
\subsection{2.2. 레슨 2-2: 평균에 대한 단측 검정 (One-Tail Test of the Means)}
\addcontentsline{toc}{subsection}{2.2. 레슨 2-2: 평균에 대한 단측 검정}

이 레슨에서는 두 모집단 평균의 차이가 특정 방향으로 나타나는지(예: 한쪽이 다른 쪽보다 크거나 작다)를 검정하는 단측 검정 방법을 학습합니다.

\subsubsection*{2.2.1. 레슨 2-2.1: 평균에 대한 단측 검정}
\addcontentsline{toc}{subsubsection}{2.2.1. 레슨 2-2.1: 평균에 대한 단측 검정}

\begin{summarybox}[title=핵심 요약: 단측 검정의 목적]
단측 검정은 두 모집단 평균의 차이가 특정 방향(예: '더 크다' 또는 '더 작다')으로 존재하는지 여부를 검정할 때 사용됩니다. 이는 양측 검정('다르다')과 달리 특정 가설 방향이 있을 때 유용합니다.
\end{summarybox}

\begin{infobox}[title=단측 검정 (One-tail Test) 가설 설정]
때로는 한 모집단의 평균이 다른 모집단의 평균과 특정 방향으로 다른지 알고 싶을 때가 있습니다. 즉, 모집단 1의 평균이 모집단 2의 평균보다 크거나 같거나, 그 반대인 경우입니다. 이런 경우 우리는 단측 검정을 수행할 것입니다. 두 가지 가능성은 다음과 같습니다:

\textbf{가능성 1: $\mu_1$이 $\mu_2$보다 작다는 것을 검정 (좌측 꼬리 검정)}
\begin{itemize}
    \item $H_0 : \mu_1 \ge \mu_2 \quad \text{또는} \quad H_0 : \mu_1 - \mu_2 \ge 0$ (귀무 가설: $\mu_1$이 $\mu_2$보다 크거나 같다)
    \item $H_A : \mu_1 < \mu_2 \quad \text{또는} \quad H_A : \mu_1 - \mu_2 < 0$ (대립 가설: $\mu_1$이 $\mu_2$보다 작다)
\end{itemize}

\textbf{가능성 2: $\mu_1$이 $\mu_2$보다 크다는 것을 검정 (우측 꼬리 검정)}
\begin{itemize}
    \item $H_0 : \mu_1 \le \mu_2 \quad \text{또는} \quad H_0 : \mu_1 - \mu_2 \le 0$ (귀무 가설: $\mu_1$이 $\mu_2$보다 작거나 같다)
    \item $H_A : \mu_1 > \mu_2 \quad \text{또는} \quad H_A : \mu_1 - \mu_2 > 0$ (대립 가설: $\mu_1$이 $\mu_2$보다 크다)
\end{itemize}
\end{infobox}

\begin{examplebox}[title=제품 테스트 (1/5): 새로운 점심 메뉴 도입]
한 패스트푸드 회사가 새로운 점심 메뉴를 고려하고 있습니다. 회사는 테스트 시장 결과를 바탕으로 메뉴를 선택할 것입니다. 이 메뉴는 기존 메뉴 항목 중 하나보다 더 높은 주간 수요를 보여야 합니다. 선택되려면 0.01 유의 수준에서 차이를 보여야 합니다.

\textbf{배경 설명:} \\
패스트푸드 체인에서 매우 흔한 시나리오를 생각해 봅시다. 새로운 메뉴 항목을 추가하기로 결정할 때마다 많은 작업과 분석이 그 결정에 들어갑니다. 새로 추가되는 모든 항목은 먼저 상당한 수요를 보여야만 추가될 수 있습니다. 새로 추가되는 모든 항목은 원자재 조달의 복잡성을 증가시키고 그에 따른 작업 흐름 재설계를 필요로 합니다. 새로운 항목은 먼저 여러 매장의 메뉴에 추가된 다음, 새로운 항목에 대한 결정이 내려지기 전에 메뉴의 다른 일반 항목의 수요와 비교됩니다. 그럼, 이러한 연구를 어떻게 수행하고 분석할지 살펴보겠습니다.
\end{examplebox}

\begin{infobox}[title=제품 테스트 (2/5): 가설 설정 (새 메뉴)]
\textbf{요구 사항:} 새로운 메뉴 항목은 기존 메뉴 항목 중 하나보다 더 잘 팔려야 합니다. 분석은 0.01 유의 수준에서 차이를 보여야 선택됩니다.

\textbf{가설:}
\begin{itemize}
    \item $H_0 : \mu_{new} \le \mu_{regular} \quad \rightarrow \quad H_0 : \mu_{new} - \mu_{regular} \le 0$ (귀무 가설: 새로운 항목의 수요가 기존 항목보다 작거나 같다)
    \item $H_A : \mu_{new} > \mu_{regular} \quad \rightarrow \quad H_A : \mu_{new} - \mu_{regular} > 0$ (대립 가설: 새로운 항목의 수요가 기존 항목보다 크다)
\end{itemize}
\textbf{유의 수준:} $\alpha = 0.01$

\textbf{설명:} \\
먼저, 이 새로운 제품이 다른 항목보다 더 잘 팔리기를 원합니다. 따라서 이것이 우리의 대립 가설이 될 것입니다. 이는 $\mu_{new} - \mu_{regular}$가 0보다 크다고 작성할 수 있습니다. 따라서 그 보완은 귀무 가설이 될 것입니다. 여기서 $\alpha$는 0.01입니다. 다시 한번, 데이터 수집은 정말 중요합니다.
\end{infobox}

\begin{infobox}[title=제품 테스트 (3/5): 데이터 수집 방법]
\textbf{데이터 수집:} 무작위로 독립적으로 수집된 데이터 (주간 판매량)
\begin{itemize}
    \item 일부 매장에는 새로운 항목이 메뉴에 있으며, 새로운 항목에 대한 수요 데이터를 제공합니다.
    \item 비교 대상인 기존 항목에 대한 수요 데이터는 아직 새로운 항목이 없는 매장에서 수집됩니다.
\end{itemize}
우리는 독립적인 두 모집단에서 표본 데이터를 수집할 것입니다. 한 모집단은 메뉴에 새로운 항목이 있고, 다른 모집단은 새로운 항목이 없습니다.
\end{infobox}

\begin{infobox}[title=제품 테스트 (4/5): 결과 분석 및 p-값 결정]
\textbf{가설:}
\begin{itemize}
    \item $H_0 : \mu_{new} - \mu_{regular} \le 0$
    \item $H_A : \mu_{new} - \mu_{regular} > 0$
\end{itemize}
\textbf{유의 수준:} $\alpha = 0.01$

\textbf{Excel 결과:} (t-검정, 두 표본, 분산 불균등 가정)
\begin{table}[h!]
    \centering
    \caption{Excel 결과: t-검정 (두 표본, 분산 불균등 가정) - 신규/기존 메뉴}
    \label{tab:excel_output_product_test}
    \begin{adjustbox}{width=\textwidth,center}
    \begin{tabular}{lcc}
        \toprule
        \textbf{} & \textbf{신규 항목 (New Item)} & \textbf{기존 항목 (Regular Item)} \\
        \midrule
        평균 (Mean) & 17.50666667 & 16.9444444 \\
        분산 (Variance) & 4.037117117 & 20.8638006 \\
        관측치 (Observations) & 75 & 108 \\
        가설 평균 차이 (Hypothesized Mean Difference) & 0 & \\
        자유도 (df) & 157 & \\
        t-통계량 (t-stat) & 2.640045854 & \\
        \textbf{P(T $\le$ t) 단측 (one-tail)} & \textbf{0.004563106} & \\
        T 임계값 단측 (critical one-tail) & 2.350333732 & \\
        P(T $\le$ t) 양측 (two-tail) & 0.009126212 & \\
        T 임계값 양측 (critical two-tail) & 2.607506474 & \\
        \bottomrule
    \end{tabular}
    \end{adjustbox}
\end{table}

\textbf{결정:}
p-값 (단측) = 0.0045 < 0.01 ($\alpha$) $\rightarrow$ 귀무 가설을 기각합니다.

\textbf{설명:} \\
데이터가 수집되면, 데이터를 기반으로 관측된 차이에 대한 p-값을 찾기 위해 분석을 수행할 것입니다. 데이터 결과는 위에 표시되어 있으며, 평균 숫자는 천 단위입니다. 또한, 새로운 항목에 대한 표본 크기가 다릅니다. 새로운 항목에 대한 주간 판매 데이터는 75개이고, 기존 항목에 대한 주간 판매 데이터는 108개입니다. 다시 한번, 각 모집단에서 독립적으로 데이터를 수집하고 있으므로 표본 크기가 동일할 필요는 없습니다.

우리는 단측 검정을 수행하고 있으므로, $\alpha$와 비교할 p-값은 0.0045입니다. 이 값은 $\alpha$ 0.01보다 작습니다. 따라서 귀무 가설을 기각할 것입니다.
\end{infobox}

\begin{infobox}[title=제품 테스트 (5/5): 최종 결론 및 경영진의 관심]
\textbf{가설:}
\begin{itemize}
    \item $H_0 : \mu_{new} - \mu_{regular} \le 0$
    \item $H_A : \mu_{new} - \mu_{regular} > 0$
\end{itemize}
\textbf{유의 수준:} $\alpha = 0.01$
\textbf{p-값:} $0.0045 < 0.01 \rightarrow$ 귀무 가설을 기각합니다 ($Reject \ H_0$).

\textbf{결론:} \\
대립 가설을 지지하여 귀무 가설을 기각하는 것은 새로운 항목이 비교 대상인 기존 항목보다 더 잘 팔린다고 믿는다는 것을 의미합니다.
따라서 경영진은 이 새로운 항목을 추가하는 것을 고려해야 합니다. 경영진은 이 소식을 듣고 기뻐하지만, 새로운 항목을 전국적으로 출시하기 전에 이 두 항목 간의 예상되는 차이가 무엇인지 알고 싶어 합니다.
\end{infobox}

\begin{infobox}[title=수요의 차이는 얼마나 될까? (1/2) - 신뢰 구간 계산 (단측 검정 후)]
\textbf{Excel 결과:} (신규/기존 메뉴 t-검정 표)

\textbf{신뢰 구간 공식:}
$\left[(\bar{x}_1 - \bar{x}_2) \pm t_{\frac{\alpha}{2}} \sqrt{\frac{s_1^2}{n_1} + \frac{s_2^2}{n_2}}\right]$

\textbf{계산 단계:}
\begin{itemize}
    \item \textbf{평균 차이 ($\bar{x}_1 - \bar{x}_2$):} $17.51 - 16.9 = 1.32$
    \item \textbf{t-임계값 (단측 검정 후 신뢰 구간):} $t_{\frac{\alpha}{2}} = 2.35$
    \textbf{설명:} 단측 검정($\alpha=0.01$)을 수행한 후 신뢰 구간을 만들 때는, 양측 꼬리 각각에 $\alpha/2$의 오류 확률을 배분하여 $1-2\alpha$ 신뢰 구간을 구성합니다. 즉, 1\% 유의 수준의 단측 검정 후에는 98\% 신뢰 구간을 만듭니다. 이때 사용되는 t-값은 Excel 출력의 'T critical one-tail' 값인 2.350333732입니다.
    \item \textbf{표준 오차 (Standard error):} $\sqrt{\frac{4.037}{75} + \frac{20.868}{108}} = 0.497$
    \item \textbf{오차 한계 (Margin of Error):} $2.35 \times 0.497 = 1.168$
\end{itemize}

\textbf{설명:} \\
우리의 결과와 두 모집단 간의 차이에 대한 신뢰 구간 방정식을 고려하면 답을 찾을 수 있습니다. 단측 검정을 수행한 후, 이제 평균 차이에 대한 신뢰 구간을 찾으려고 합니다. 이때 구간의 하한과 상한을 모두 추정하는 데 오류가 있을 수 있으므로, 단측 t-값(이 경우 0.01)을 사용하면 양쪽 꼬리 각각에서 1\%의 오류 확률을 갖게 됩니다. 따라서 이 경우 신뢰 구간은 98\% 신뢰 구간이 됩니다. 여기서는 이 차이를 단계별로 계산하는 방법을 보여줍니다.

첫 번째는 두 표본 평균의 차이입니다. 다음으로 오차 한계가 필요합니다. 이는 단측 검정의 $t_{\frac{\alpha}{2}}$ 값(우리 경우 2.35)에 표준 오차를 곱한 것입니다. 이 경우 오차 한계는 1.168이 됩니다. 이제 98\% 신뢰 구간을 표현할 준비가 되었습니다. 여기서 $\alpha$는 0.01입니다.
\end{infobox}

\begin{infobox}[title=수요의 차이는 얼마나 될까? (2/2) - 최종 신뢰 구간 및 해석]
\textbf{계산된 값:}
\begin{itemize}
    \item 평균 차이 ($\bar{x}_1 - \bar{x}_2$): $17.51 - 16.9 = 1.32$
    \item 오차 한계 (Margin of Error): $2.35 \times 0.497 = 1.168$
\end{itemize}
\textbf{98\% 신뢰 구간:} $1.32 \pm 1.168 = [0.1517, 2.488]$

\textbf{천 단위로 재작성:}
주간 판매량이 천 단위로 기록되었으므로, 이를 다음과 같이 다시 작성할 수 있습니다: $[151.7, 2,488.22]$

\textbf{해석:} \\
우리는 98\% 신뢰 구간이 0.1517에서 2.488 사이임을 발견했습니다. 주간 판매량이 천 단위로 기록되었으므로, 이를 151.7에서 2,488.22로 다시 작성할 수 있습니다. 즉, 우리는 새로운 항목의 주간 평균 수요가 기존 항목보다 최소 151.7단위 더 많다고 98\% 확신할 수 있습니다.

이 구간의 폭이 얼마나 넓은지 보세요. 이 구간의 폭이 마음에 들지 않을 수도 있습니다. 이는 부분적으로 우리가 요구한 신뢰 수준 때문입니다. 신뢰 수준이 높을수록 이 구간은 더 넓어진다는 것을 기억하세요. 하지만 대체로 우리는 이 새로운 항목이 비교 대상인 기존 항목보다 더 잘 팔릴 것으로 예상합니다.
\end{infobox}

\begin{examplebox}[title=연습 문제: 교육 프로그램 효과 - 문제 상황]
한 교육 프로그램은 대부분의 직원의 판매 실적을 향상시킬 것을 약속합니다. 영업 부서 관리자는 직원 80명을 이 프로그램에 참여시킵니다. 무작위로 선정된 대조군보다 5\% 유의 수준에서 결과가 더 좋다고 나타나면, 모든 직원을 프로그램에 참여시킬 것입니다.

\textbf{목표:} 이 문제에 대한 귀무 가설과 대립 가설을 설정하세요.
\begin{itemize}
    \item 그룹 1: 대조군 ($\mu_1$: 대조군의 주간 평균 판매액)
    \item 그룹 2: 교육받은 직원 ($\mu_2$: 교육을 받은 직원의 주간 평균 판매액)
\end{itemize}
\end{examplebox}

\begin{infobox}[title=연습 문제: 귀무 가설 예시 (선택)]
한 교육 프로그램은 대부분의 직원의 판매 실적을 향상시킬 것을 약속합니다. 영업 부서 관리자는 직원 80명을 이 프로그램에 참여시킵니다. 무작위로 선정된 대조군보다 5\% 유의 수준에서 결과가 더 좋다고 나타나면, 모든 직원을 프로그램에 참여시킬 것입니다.

\textbf{목표:} 이 문제에 대한 귀무 가설과 대립 가설을 설정하세요.
\begin{itemize}
    \item $\mu_1$: 대조군의 주간 평균 판매액
    \item $\mu_2$: 교육을 받은 직원의 주간 평균 판매액
\end{itemize}

\textbf{선택지:}
\begin{enumerate}
    \item $H_0 : \mu_1 - \mu_2 = 0$ 그리고 $H_A : \mu_1 - \mu_2 \neq 0$
    \item \textbf{$H_0 : \mu_1 - \mu_2 \ge 0$ 그리고 $H_A : \mu_1 - \mu_2 < 0$}
    \item $H_0 : \mu_1 - \mu_2 \le 0$
\end{enumerate}
\end{infobox}

\begin{infobox}[title=연습 문제: 풀이 - 가설 설정 (교육 프로그램)]
교육 프로그램은 대부분의 직원의 판매 실적을 향상시킬 것을 약속하므로, 우리는 데이터가 $\mu_1$이 $\mu_2$보다 작다는 것을 보여주기를 원합니다. 즉, 교육받은 직원의 판매 실적이 대조군보다 높기를 기대합니다. 따라서 이것이 우리의 대립 가설이 될 것이고, 그 보완은 귀무 가설이 될 것입니다.

\textbf{가설:}
\begin{itemize}
    \item $H_0 : \mu_1 \ge \mu_2 \quad \rightarrow \quad H_0 : \mu_1 - \mu_2 \ge 0$ (귀무 가설: 대조군의 판매 실적이 교육받은 직원보다 크거나 같다)
    \item $H_A : \mu_1 < \mu_2 \quad \rightarrow \quad H_A : \mu_1 - \mu_2 < 0$ (대립 가설: 대조군의 판매 실적이 교육받은 직원보다 작다)
\end{itemize}
\end{infobox}

\begin{infobox}[title=연습 문제: 결과 분석 (교육 프로그램)]
\textbf{가설:}
\begin{itemize}
    \item $H_0 : \mu_1 \ge \mu_2$
    \item $H_A : \mu_1 < \mu_2$
\end{itemize}

\textbf{교육 프로그램 결과:}
\begin{table}[h!]
    \centering
    \caption{교육 프로그램 결과: t-검정 (두 표본, 분산 불균등 가정)}
    \label{tab:excel_output_training_program}
    \begin{adjustbox}{width=\textwidth,center}
    \begin{tabular}{lcc}
        \toprule
        \textbf{} & \textbf{그룹 1 (대조군)} & \textbf{그룹 2 (새로 교육받은 직원)} \\
        \midrule
        평균 (Mean) & 1419.580645 & 1664.575 \\
        분산 (Variance) & 69082.15919 & 67313.69051 \\
        관측치 (Observations) & 93 & 80 \\
        가설 평균 차이 (Hypothesized Mean Difference) & 0 & \\
        자유도 (df) & 168 & \\
        t-통계량 (t-stat) & -6.155248397 & \\
        \textbf{P(T $\le$ t) 단측 (one-tail)} & \textbf{2.66462E-09} & \\
        T 임계값 단측 (critical one-tail) & 1.653974208 & \\
        P(T $\le$ t) 양측 (two-tail) & 5.32925E-09 & \\
        T 임계값 양측 (critical two-tail) & 1.974185191 & \\
        \bottomrule
    \end{tabular}
    \end{adjustbox}
\end{table}

\textbf{질문:} 5\% 유의 수준에서 당신의 결론은 무엇입니까?
\end{infobox}

\begin{infobox}[title=연습 문제: 예시 표 (선택)]
데이터가 분석되어 여기에 제시되어 있습니다. 5\% 유의 수준에서 당신의 결론은 무엇입니까? (슬라이드 55의 표 데이터 참조)
\begin{enumerate}
    \item \textbf{귀무 가설을 기각합니다 (Reject the null hypothesis)}
    \item 귀무 가설을 기각하지 않습니다 (Not reject the null hypothesis)
\end{enumerate}
\end{infobox}

\begin{infobox}[title=연습 문제: 풀이 - 결론 (교육 프로그램)]
\textbf{가설:}
\begin{itemize}
    \item $H_0 : \mu_1 \ge \mu_2$
    \item $H_A : \mu_1 < \mu_2$
\end{itemize}
\textbf{결정:}
P(T $\le$ t) 단측 = 2.66462E-09 < 0.05 ($\alpha$) $\rightarrow$ 귀무 가설을 기각합니다.

\textbf{설명:} \\
이것은 단측 검정입니다. 5\% 유의 수준에서 이와 같은 결과를 볼 p-값은 매우 작으며, 이는 대립 가설을 지지하여 귀무 가설을 기각하게 할 것입니다.
교육이 영업 사원들이 더 효과적이 되도록 도왔다는 것이 분명해 보입니다. 관리자는 모든 직원을 위한 교육 프로그램을 고려해야 합니다.
\end{infobox}

\begin{examplebox}[title=연습 문제: 신뢰 구간 계산 (교육 프로그램)]
교육 프로그램은 많은 비용이 듭니다. 이 프로그램에 전념하기로 결정하기 전에, 관리자는 이 교육을 통해 예상할 수 있는 주간 판매액의 차이를 알고 싶어 합니다. 차이에 대한 90\% 신뢰 구간을 개발하세요.

\textbf{참고:} 신뢰 구간의 경우 5\%는 양쪽 꼬리에 모두 있을 것이므로, 여기서 추적하는 신뢰 구간은 90\% 신뢰 구간이 될 것입니다. 90\% 신뢰 구간을 계산하기 위해 이 출력을 사용하세요.
\end{examplebox}

\begin{infobox}[title=연습 문제: 풀이 - 신뢰 구간 (교육 프로그램)]
\textbf{교육 프로그램 결과:} (슬라이드 55의 표 데이터 참조)

\textbf{계산 단계:}
\begin{itemize}
    \item \textbf{평균 차이 ($\bar{x}_1 - \bar{x}_2$):} $1419.58 - 1664.58 = -245$
    \item \textbf{t-임계값 (단측 검정 후 신뢰 구간):} $t_{\frac{\alpha}{2}} = 1.654$
    \textbf{설명:} 90\% 신뢰 구간이므로, 양측 꼬리 각각 5%에 해당하는 t-값을 사용합니다. Excel 출력의 'T critical one-tail' 값 1.653974208을 사용합니다.
    \item \textbf{표준 오차 (Standard Error):} $\sqrt{\frac{69082.16}{93} + \frac{67313.69}{80}} = 39.8$
    \item \textbf{오차 한계 (Margin of Error):} $1.654 \times 39.8 = 65.83$
\end{itemize}
\textbf{90\% 신뢰 구간:} $[-245 \pm 65.83] = [-310.83, -179.17]$

\textbf{설명:} \\
다시 한번, 신뢰 구간의 경우 5\%는 양쪽 꼬리에 모두 있을 것이므로, 여기서 계산하는 신뢰 구간은 90\%입니다. 90\% 신뢰 구간을 계산하기 위해 출력을 사용하면, 먼저 이 두 표본 간의 차이는 -245입니다. 다음으로 $t_{\frac{\alpha}{2}}$는 1.654이고, 표준 오차는 39.8입니다. 이제 오차 한계를 65.83으로 계산할 수 있으며, 이는 90\% 신뢰 구간을 -310.83에서 -179.17로 만듭니다.
\end{infobox}

\begin{infobox}[title=무엇을 의미하는가? (What Does It Mean?) - 교육 프로그램 신뢰 구간 해석]
\textbf{90\% 신뢰 구간:} $[-310.83, -179.17]$

\textbf{해석:}
\begin{itemize}
    \item 교육을 받지 않은 직원들은 평균적으로 판매량이 더 낮습니다 (달러 기준).
    \item 실제 차이 (90\% 신뢰 수준)는 교육을 받은 그룹에 비해 $179.17$달러에서 $310.83$달러 더 적습니다. (즉, 교육받은 그룹이 대조군보다 평균적으로 이만큼 더 많이 판매합니다.)
\end{itemize}

\textbf{설명:} \\
이것이 정확히 무엇을 의미할까요? 이것은 첫 번째 그룹인 대조군과 교육받은 그룹 간의 평균 차이라는 것을 기억하세요. 대조군은 평균적으로 교육받은 그룹보다 판매량이 낮았습니다. 그래서 우리는 교육을 고려하고 있는 것입니다. 하지만 경영진은 이 프로그램에 자원을 투입하기 전에 차이가 어느 정도일지 알고 싶어 했습니다. 이 경우, 대조군의 평균 판매액은 교육받은 그룹보다 적으며, 90\% 신뢰 수준에서 실제 값은 $179.17$달러에서 $310.83$달러 사이입니다.

만약 관리자가 이것이 좋기는 하지만 비용을 정당화할 만큼 충분히 좋지 않다고 생각한다면 어떨까요? 특정 수준의 차이를 찾는 가설을 이런 방식으로 하는 대신 검정할 수 있을까요? 답은 '예'입니다. 그리고 그것이 우리가 다음 레슨에서 배울 내용입니다.
\end{infobox}

\newpage

\newpage
\section{개요}
이 문서는 UIUC BADM-572 과목의 '추론 및 예측 통계 모듈 2' 중 파트 2/3에 해당하는 내용을 다룹니다.
주요 내용은 평균에 대한 단측 검정, 특정 값에 대한 평균 검정, 짝지은 표본 검정, 그리고 비율 검정입니다.
각 개념은 비전공자도 쉽게 이해할 수 있도록 상세한 설명, 현실적인 예시, 그리고 Excel을 활용한 실습 절차를 포함합니다.
가설 설정부터 결과 해석, 그리고 신뢰 구간 계산까지 전반적인 통계적 의사결정 과정을 학습하는 데 중점을 둡니다.

\newpage
\section{용어 정리}
\begin{adjustbox}{width=\textwidth,center}
\begin{tabular}{llll}
\toprule
\textbf{용어} & \textbf{쉬운 설명} & \textbf{원어} & \textbf{비고} \\
\midrule
단측 검정 & 한 방향으로의 차이만 검정하는 가설 검정 (예: 더 크다, 더 작다) & One-Tail Test & 특정 방향의 변화에 관심 있을 때 사용 \\
양측 검정 & 양방향으로의 차이를 검정하는 가설 검정 (예: 다르다) & Two-Tail Test & 차이의 방향에 관계없이 변화 여부만 볼 때 사용 \\
유의 수준 ($\alpha$) & 귀무 가설을 잘못 기각할 확률의 최대 허용치 & Level of Significance & 일반적으로 0.05 또는 0.01 사용 \\
P-값 & 귀무 가설이 참일 때, 관측된 데이터 또는 더 극단적인 데이터가 나올 확률 & P-value & P-값이 $\alpha$보다 작으면 귀무 가설 기각 \\
귀무 가설 ($H_0$) & 통계적 검정의 기본 가정 (차이가 없다, 효과가 없다 등) & Null Hypothesis & 기각되지 않으면 현상 유지 \\
대립 가설 ($H_A$) & 귀무 가설과 반대되는 주장 (차이가 있다, 효과가 있다 등) & Alternative Hypothesis & 연구자가 증명하고자 하는 가설 \\
t-통계량 & 표본 평균이 귀무 가설의 평균과 얼마나 떨어져 있는지 나타내는 값 & t-statistic & t-분포를 따르며, P-값 계산에 사용 \\
t-임계값 & 유의 수준 $\alpha$에서 귀무 가설을 기각하기 위한 t-통계량의 기준 값 & T critical value & t-통계량이 t-임계값을 넘어서면 귀무 가설 기각 \\
자유도 (df) & 통계량을 계산할 때 자유롭게 변할 수 있는 데이터의 수 & Degrees of Freedom & 표본 크기에 따라 달라짐 \\
표준 오차 & 표본 통계량(예: 표본 평균)이 모수(예: 모평균)와 얼마나 차이 날 수 있는지 나타내는 척도 & Standard Error & 표본 변동성의 추정치 \\
신뢰 구간 & 모수가 특정 범위 내에 있을 것이라고 추정되는 구간 & Confidence Interval & 신뢰 수준(예: 95\%)으로 표현 \\
짝지은 표본 검정 & 동일한 개체에 대해 두 번 측정하거나, 특성이 유사한 두 개체를 짝지어 비교하는 검정 & Paired Test & 개체 간 변동성을 제거하여 비교 정확도 높임 \\
비율 검정 & 두 모집단의 특정 특성(예: 성공률, 비율)이 서로 다른지 검정 & Test of Proportions & Z-검정을 주로 사용 \\
\bottomrule
\end{tabular}
\end{adjustbox}

\newpage
\section{Lesson 2-2.2: Excel을 이용한 평균의 단측 검정}

\subsection{개념: 단측 검정의 필요성}
우리는 특정 방향으로의 변화에만 관심이 있을 때 단측 검정을 사용합니다. 예를 들어, 새로운 제품이 기존 제품보다 '더 높은' 수요를 보일지, 또는 새로운 치료법이 기존 치료법보다 '더 빠른' 회복 시간을 가져올지 등을 알고 싶을 때 유용합니다. 이는 단순히 '다르다'는 것을 넘어 '어떤 방향으로 다른지'를 확인하는 데 초점을 맞춥니다.

\begin{examplebox}[title=시나리오: 패스트푸드 신메뉴 출시]
한 패스트푸드 회사가 신메뉴 출시를 고려하고 있습니다. 이 신메뉴는 기존 메뉴 중 하나보다 '더 높은' 주간 수요를 보여야만 선택될 것입니다. 분석은 유의 수준 0.01에서 차이를 보여야 합니다.
\end{examplebox}

\subsection{가설 설정}
가설 검정의 첫 단계는 귀무 가설($H_0$)과 대립 가설($H_A$)을 설정하는 것입니다.

\begin{itemize}
    \item \textbf{귀무 가설 ($H_0$):} 신메뉴의 평균 수요가 기존 메뉴의 평균 수요보다 작거나 같다.
    ($\mu_{New} \le \mu_{Current} \implies \mu_{New} - \mu_{Current} \le 0$)
    \item \textbf{대립 가설 ($H_A$):} 신메뉴의 평균 수요가 기존 메뉴의 평균 수요보다 높다.
    ($\mu_{New} > \mu_{Current} \implies \mu_{New} - \mu_{Current} > 0$)
\end{itemize}
우리는 신메뉴가 '더 높은' 수요를 보이는지 확인하고 싶으므로, 이는 단측 검정(one-tail test)이 됩니다.

\subsection{Excel을 이용한 분석 절차}
Excel의 '데이터 분석' 도구를 사용하여 두 표본 평균의 차이에 대한 t-검정을 수행할 수 있습니다.

\begin{enumerate}
    \item \textbf{데이터 준비:} 신메뉴와 기존 메뉴의 주간 수요 데이터를 Excel 워크시트에 준비합니다. (예: 신메뉴 데이터는 A열, 기존 메뉴 데이터는 B열)
    \item \textbf{데이터 분석 도구 실행:}
    \begin{itemize}
        \item '데이터' 탭으로 이동합니다.
        \item '데이터 분석'을 선택합니다.
        \item 't-검정: 등분산 가정 두 표본' 또는 't-검정: 이분산 가정 두 표본'을 선택합니다. (여기서는 't-검정: 이분산 가정 두 표본'을 사용합니다.)
    \end{itemize}
    \item \textbf{입력 설정:}
    \begin{itemize}
        \item \textbf{변수 1 범위 (Variable 1 Range):} 신메뉴 데이터가 있는 열을 선택합니다 (예: \$A\$A).
        \item \textbf{변수 2 범위 (Variable 2 Range):} 기존 메뉴 데이터가 있는 열을 선택합니다 (예: \$B\$B).
        \item \textbf{가설 평균 차이 (Hypothesized Mean Difference):} 0을 입력합니다. (귀무 가설에서 두 평균의 차이가 0이거나 0보다 작다고 가정하기 때문입니다.)
        \item \textbf{레이블 (Labels):} 데이터에 열 제목이 포함되어 있다면 이 옵션을 선택합니다.
        \item \textbf{알파 (Alpha):} 유의 수준인 0.01을 입력합니다.
        \item \textbf{출력 옵션 (Output Options):} 결과를 표시할 위치를 선택합니다 (예: '새 워크시트' 또는 '출력 범위').
    \end{itemize}
    \item \textbf{결과 확인:} '확인'을 클릭하여 분석 결과를 얻습니다.
\end{enumerate}

\subsection{결과 해석}
Excel 출력 결과는 다음과 같은 정보를 포함합니다.

\begin{table}[h!]
\caption{품목 판매 결과 (t-검정: 이분산 가정 두 표본)}
\label{tab:fastfood_ttest}
\centering
\begin{adjustbox}{width=\textwidth,center}
\begin{tabular}{lcc}
\toprule
\textbf{항목} & \textbf{신메뉴} & \textbf{기존 메뉴} \\
\midrule
평균 (Mean) & 17.50666667 & 16.19444444 \\
분산 (Variance) & 4.037117117 & 620.86838006 \\
관측치 (Observations) & 75 & 108 \\
가설 평균 차이 (Hypothesized Mean Difference) & 0 & \\
자유도 (df) & 157 & \\
t-통계량 (t-stat) & 2.640045854 & \\
P(T $\le$ t) 단측 (P(T $\le$ t) one-tail) & \textbf{0.004563106} & \\
T 임계값 단측 (T critical one-tail) & \textbf{2.350333732} & \\
P(T $\le$ t) 양측 (P(T $\le$ t) two-tail) & 0.009126212 & \\
T 임계값 양측 (T critical two-tail) & 2.607506474 & \\
\bottomrule
\end{tabular}
\end{adjustbox}
\end{table}

\begin{itemize}
    \item \textbf{P-값 (단측):} 우리의 대립 가설은 신메뉴의 수요가 '더 높다'는 단측 가설이므로, 'P(T $\le$ t) 단측' 값을 확인합니다. 이 값은 0.00456입니다.
    \item \textbf{유의 수준 ($\alpha$):} 우리는 0.01의 유의 수준을 설정했습니다.
    \item \textbf{결론:} P-값 (0.00456)이 유의 수준 $\alpha$ (0.01)보다 작습니다.
    ($0.00456 < 0.01$)
\end{itemize}
따라서, 우리는 귀무 가설을 기각합니다. 이는 신메뉴의 수요가 기존 메뉴보다 통계적으로 유의미하게 더 높다는 것을 의미합니다.

\subsection{신뢰 구간 계산}
신메뉴가 기존 메뉴보다 얼마나 더 높은 수요를 보이는지 추정하기 위해 신뢰 구간을 계산할 수 있습니다.

\begin{enumerate}
    \item \textbf{평균 차이 ($\bar{d}$):}
    신메뉴 평균 - 기존 메뉴 평균 = $17.50666667 - 16.19444444 = 1.312222222$ (천 단위)
    \item \textbf{t-값 ($t_{\alpha/2}$):}
    단측 검정에서 $\alpha = 0.01$을 사용했으므로, 양측 신뢰 구간을 계산할 때는 양쪽 꼬리에 각각 0.01씩, 즉 총 0.02를 제외한 98\% 신뢰 구간을 얻게 됩니다. 이때 사용되는 t-값은 'T 임계값 단측'인 2.350333732입니다.
    \item \textbf{표준 오차 (Standard Error):}
    표준 오차는 다음 공식으로 계산됩니다:
    $\text{표준 오차} = \sqrt{\frac{\text{신메뉴 분산}}{\text{신메뉴 관측치}} + \frac{\text{기존 메뉴 분산}}{\text{기존 메뉴 관측치}}}$
    $\text{표준 오차} = \sqrt{\frac{4.037117117}{75} + \frac{620.86838006}{108}} = \sqrt{0.053828228 + 5.748781297} = \sqrt{5.802609525} \approx 0.497045239$
    \item \textbf{오차 한계 (Margin of Error):}
    오차 한계 = t-값 $\times$ 표준 오차 = $2.350333732 \times 0.497045239 \approx 1.168222192$
    \item \textbf{신뢰 구간:}
    하한 (Lower Bound) = 평균 차이 - 오차 한계 = $1.312222222 - 1.168222192 = 0.14400003$
    상한 (Upper Bound) = 평균 차이 + 오차 한계 = $1.312222222 + 1.168222192 = 2.480444414$
\end{enumerate}

\begin{summarybox}[title=98\% 신뢰 구간 해석]
우리는 신메뉴의 주간 수요가 기존 메뉴보다 평균적으로 0.144천 개(144개)에서 2.480천 개(2,480개) 더 높을 것이라고 98\% 신뢰합니다.
\end{summarybox}

\begin{warningbox}[title=경영진의 의사결정]
신뢰 구간이 [144, 2480]으로 상당히 넓습니다. 이는 신메뉴가 최소 144개 더 팔릴 수도 있지만, 최대 2480개까지 더 팔릴 수도 있다는 의미입니다. 신메뉴 도입에는 많은 비용이 수반되므로, 경영진은 이처럼 넓은 범위의 예측에 대해 불편함을 느낄 수 있습니다. 더 정확한 정보를 얻기 위해서는 표본 크기를 늘려야 할 필요가 있습니다.
\end{warningbox}

\newpage
\section{Lesson 2-3.1: 평균의 특정 값 검정}

\subsection{개념: 특정 값과의 차이 검정}
이전까지는 두 모집단 평균의 차이가 '0'인지 여부를 검정했습니다. 하지만 때로는 두 평균의 차이가 '특정 값' 이상 또는 이하인지, 혹은 '특정 값'과 같은지 알고 싶을 때가 있습니다. 예를 들어, 새로운 공정이 기존 공정보다 생산량을 '10개 이상' 늘리는지, 또는 새로운 서비스가 고객 대기 시간을 '5분 이하'로 줄이는지 등을 검정할 수 있습니다. 이는 단순히 차이의 유무를 넘어, 특정 목표치 달성 여부를 판단하는 데 중요합니다.

\begin{examplebox}[title=시나리오: 유지보수 공급업체 변경 고려]
한 회사가 유지보수 공급업체와의 계약 만료를 앞두고 다른 회사를 고려 중입니다. 현재 공급업체(공급업체 A)의 문제 해결 시간에 대한 불만이 있습니다. 새로운 공급업체(공급업체 B)는 연간 20만 달러의 추가 비용이 들지만, 더 빠른 응답 시간을 약속합니다. 경영진은 공급업체 B로 전환하기 위해 최소 5시간의 수리 시간 단축을 원합니다. 회사는 공급업체를 변경해야 할까요?
\end{examplebox}

\subsection{가설 설정 (단계 1)}
경영진은 공급업체 B의 추가 비용을 정당화하기 위해 수리 시간이 최소 5시간 단축되기를 원합니다.
\begin{itemize}
    \item $\mu_A$: 공급업체 A의 평균 문제 해결 시간
    \item $\mu_B$: 공급업체 B의 평균 문제 해결 시간
\end{itemize}
공급업체 B가 A보다 5시간 이상 빠르다는 것은 $\mu_B \le \mu_A - 5$ 또는 $\mu_A - \mu_B \ge 5$를 의미합니다.
우리는 공급업체 B가 '충분히' 빠르다는 것을 증명해야 하므로, 대립 가설은 우리가 증명하고자 하는 주장이 됩니다.

\begin{itemize}
    \item \textbf{귀무 가설 ($H_0$):} 공급업체 A와 B의 평균 문제 해결 시간 차이가 5시간 이상이다.
    ($\mu_A - \mu_B \ge 5$)
    \item \textbf{대립 가설 ($H_A$):} 공급업체 A와 B의 평균 문제 해결 시간 차이가 5시간 미만이다.
    ($\mu_A - \mu_B < 5$)
\end{itemize}
이것은 공급업체 B가 A보다 5시간 미만으로 빠르다는 것을 증명하는 단측 검정입니다. 즉, 5시간 단축 목표를 달성하지 못했다는 것을 보여주는 것입니다.

\subsection{유의 수준 설정 (단계 2)}
경영진은 5\% 유의 수준($\alpha = 0.05$)에서 이 주장을 확인하고자 합니다.

\subsection{데이터 수집 (단계 3 준비)}
가설 검정을 진행하기 위해서는 두 공급업체로부터 문제 해결 시간에 대한 대표성 있는 표본 데이터를 무작위로 독립적으로 수집해야 합니다. 현재 공급업체의 과거 기록과 새로운 공급업체의 유사 고객 데이터를 활용할 수 있습니다.

\subsection{P-값 계산 (단계 3)}
P-값은 관측된 표본 결과 또는 더 극단적인 결과가 귀무 가설이 참일 때 나타날 확률입니다. 이는 t-통계량을 통해 계산됩니다.

\begin{itemize}
    \item \textbf{t-통계량 공식:}
    $t = \frac{(\bar{x}_1 - \bar{x}_2) - D_0}{\sqrt{\frac{s_1^2}{n_1} + \frac{s_2^2}{n_2}}}$
    여기서,
    \begin{itemize}
        \item $\bar{x}_1, \bar{x}_2$: 두 표본의 평균
        \item $D_0$: 가설 평균 차이 (이 예시에서는 5)
        \item $s_1^2, s_2^2$: 두 표본의 분산
        \item $n_1, n_2$: 두 표본의 크기
    \end{itemize}
    \item \textbf{Excel 활용:} Excel의 '데이터 분석' 도구에서 't-검정: 이분산 가정 두 표본'을 선택하고, '가설 평균 차이'에 5를 입력하여 P-값을 얻습니다.
\end{itemize}

\subsection{Excel 결과 및 해석}
Excel 분석 결과는 다음과 같습니다.

\begin{table}[h!]
\caption{공급업체 A/B t-검정: 이분산 가정 두 표본}
\label{tab:supplier_ttest}
\centering
\begin{adjustbox}{width=\textwidth,center}
\begin{tabular}{lcc}
\toprule
\textbf{항목} & \textbf{공급업체 A (시간)} & \textbf{공급업체 B (시간)} \\
\midrule
평균 (Mean) & 23.422222222 & 22.025 \\
분산 (Variance) & 21.439801 & 20.29348739 \\
관측치 (Observations) & 135 & 120 \\
가설 평균 차이 (Hypothesized Mean Difference) & \textbf{5} & \\
자유도 (df) & 251 & \\
t-통계량 (t-stat) & -6.291435152 & \\
P(T $\le$ t) 단측 (P(T $\le$ t) one-tail) & \textbf{6.95E-10} & \\
T 임계값 단측 (T critical one-tail) & 1.650947025 & \\
P(T $\le$ t) 양측 (P(T $\le$ t) two-tail) & 1.39019E-09 & \\
T 임계값 양측 (T critical two-tail) & 1.969460227 & \\
\bottomrule
\end{tabular}
\end{adjustbox}
\end{table}

\begin{itemize}
    \item \textbf{가설 평균 차이:} Excel 입력 시 5로 설정되었습니다.
    \item \textbf{표본 평균:} 공급업체 A는 23.42시간, 공급업체 B는 22.05시간입니다. 공급업체 B가 더 빠르지만, 충분히 빠른지는 P-값을 통해 판단합니다.
    \item \textbf{t-통계량:} -6.29입니다.
    \item \textbf{P-값 (단측):} $6.95 \times 10^{-10}$ (거의 0에 가까운 값)
\end{itemize}

\subsection{결정 (단계 4)}
\begin{itemize}
    \item \textbf{귀무 가설 ($H_0$):} $\mu_A - \mu_B \ge 5$
    \item \textbf{대립 가설 ($H_A$):} $\mu_A - \mu_B < 5$
    \item \textbf{유의 수준 ($\alpha$):} 0.05
    \item \textbf{P-값:} $0.000000000695097$
\end{itemize}
P-값 ($6.95 \times 10^{-10}$)이 유의 수준 $\alpha$ (0.05)보다 훨씬 작습니다.
($6.95 \times 10^{-10} < 0.05$)

\begin{summarybox}[title=결론]
귀무 가설을 기각합니다. 이는 공급업체 B가 수리 시간을 최소 5시간 단축할 수 있음을 보여주지 못했음을 의미합니다. 따라서 추가 비용을 정당화할 만한 충분한 근거가 없으므로, 회사는 공급업체 B로 전환하지 않아야 합니다.
\end{summarybox}

\newpage
\section{Lesson 2-4.1: 짝지은 표본 검정}

\subsection{개념: 짝지은 표본 실험의 이해}
지금까지는 서로 독립적인 두 모집단에서 추출된 표본을 비교하는 방법을 배웠습니다. 하지만 실제 상황에서는 표본들이 서로 '짝지어져' 있거나 '의존적'인 경우가 많습니다. 짝지은 표본 검정(Paired Test)은 이러한 상황에서 두 처리(treatment) 또는 두 시점(before/after) 간의 차이를 분석하는 데 사용됩니다. 이는 개체 간의 자연적인 변동성을 제거하여, 우리가 관심 있는 처리 효과나 변화를 더 정확하게 측정할 수 있게 해줍니다.

\begin{examplebox}[title=짝지은 표본 실험의 세 가지 시나리오]
\begin{enumerate}
    \item \textbf{하나의 표본, 두 가지 처리:} 동일한 개체들에게 두 가지 다른 처리를 적용하고 결과를 비교합니다.
    \begin{itemize}
        \item \textbf{예시:} 피실험자들에게 두 가지 다른 제품을 사용하게 한 후 각각의 만족도를 평가하도록 합니다. (예: A 제품 사용 후 점수, B 제품 사용 후 점수)
    \end{itemize}
    \item \textbf{하나의 표본, 하나의 처리, 전/후 비교:} 동일한 개체들에게 하나의 처리를 적용한 후, 처리 전과 후의 결과를 비교합니다.
    \begin{itemize}
        \item \textbf{예시:} 직원들에게 특정 교육을 받게 한 후, 교육 전과 후의 생산성을 비교합니다. (예: 교육 전 생산성, 교육 후 생산성)
    \end{itemize}
    \item \textbf{두 개의 짝지은 표본:} 특정 특성(예: 나이, 능력)이 유사한 개체들을 짝지어 두 그룹으로 나눈 후, 각 그룹에 다른 처리를 적용하고 결과를 비교합니다.
    \begin{itemize}
        \item \textbf{예시:} 독해 능력이 비슷한 아이들을 짝지어 두 그룹으로 나눈 후, 각 그룹에 다른 독해 교육 방법을 적용하고 독해력 향상 정도를 비교합니다.
    \end{itemize}
\end{enumerate}
이러한 시나리오들에서 두 데이터 표본은 서로 독립적으로 개발되지 않았으므로, 독립 표본 검정과는 다른 분석 방법이 필요합니다.
\end{examplebox}

\subsection{짝지은 표본 검정의 가설}
짝지은 표본 검정에서는 각 쌍의 '차이(difference)'에 초점을 맞춥니다. 각 참가자(또는 쌍)에 대해 두 변수($x_1, x_2$)의 차이 $d = x_1 - x_2$를 계산하고, 이 차이들의 평균($\mu_d$)에 대한 가설을 설정합니다.

\begin{table}[h!]
\caption{짝지은 표본 검정의 가설 유형}
\label{tab:paired_hypotheses}
\centering
\begin{adjustbox}{width=0.8\textwidth,center}
\begin{tabular}{lll}
\toprule
\textbf{$H_0$} & \textbf{$H_A$} & \textbf{유형} \\
\midrule
$\mu_d = D$ & $\mu_d \ne D$ & 양측 검정 (Two-tail) \\
$\mu_d \ge D$ & $\mu_d < D$ & 단측 검정 (One-tail) \\
$\mu_d \le D$ & $\mu_d > D$ & 단측 검정 (One-tail) \\
\bottomrule
\end{tabular}
\end{adjustbox}
\end{table}
여기서 $D$는 가설에서 설정하는 차이의 특정 값입니다. 대부분의 경우 $D=0$으로 설정하여 차이가 없는지 여부를 검정합니다.

\begin{examplebox}[title=시나리오: 제품 디자인 평가]
한 회사가 두 가지 제품 디자인(A와 B)을 고려 중입니다. 전국 포커스 그룹 참가자들은 두 제품을 여러 차원에서 평가하고 점수(100점 만점)를 매겼습니다. 유의 수준 1\%에서 한 디자인이 다른 디자인보다 우수한지 분석합니다.
\end{examplebox}

\subsection{제품 디자인 평가: 가설 설정}
우리는 어떤 디자인이 더 우수할지에 대한 사전 가정이 없으므로, 두 디자인의 평가 점수가 동일할 것이라는 귀무 가설을 설정합니다.
\begin{itemize}
    \item \textbf{귀무 가설 ($H_0$):} 두 제품 디자인의 평균 평가 점수 차이($\mu_d$)는 0이다. ($\mu_d = 0$)
    \item \textbf{대립 가설 ($H_A$):} 두 제품 디자인의 평균 평가 점수 차이($\mu_d$)는 0이 아니다. ($\mu_d \ne 0$)
    \item \textbf{유의 수준 ($\alpha$):} 0.05 (슬라이드에는 0.01로 언급되었으나, 이후 설명에서 0.05로 진행되므로 0.05를 따릅니다.)
\end{itemize}
이는 양측 검정(two-tail test)입니다.

\subsection{Excel 결과 및 해석}
Excel의 't-검정: 짝지은 두 표본 평균' 도구를 사용하여 분석한 결과는 다음과 같습니다.

\begin{table}[h!]
\caption{제품 A/B t-검정: 짝지은 두 표본 평균}
\label{tab:product_ttest_paired}
\centering
\begin{adjustbox}{width=\textwidth,center}
\begin{tabular}{lcc}
\toprule
\textbf{항목} & \textbf{제품 A} & \textbf{제품 B} \\
\midrule
평균 (Mean) & 61.752 & 57.312 \\
분산 (Variance) & 538.1310201 & 542.9785703 \\
관측치 (Observations) & 250 & 250 \\
Pearson 상관 계수 (Pearson Correlation) & -0.063074728 & \\
가설 평균 차이 (Hypothesized Mean Difference) & 0 & \\
자유도 (df) & 249 & \\
t-통계량 (t-stat) & 2.070791373 & \\
P(T $\le$ t) 단측 (P(T $\le$ t) one-tail) & 0.019704445 & \\
T 임계값 단측 (T critical one-tail) & 1.6509996152 & \\
P(T $\le$ t) 양측 (P(T $\le$ t) two-tail) & \textbf{0.03940889} & \\
T 임계값 양측 (T critical two-tail) & 1.969536868 & \\
\bottomrule
\end{tabular}
\end{adjustbox}
\end{table}

\begin{itemize}
    \item \textbf{P-값 (양측):} 우리의 가설은 양측 검정이므로, 'P(T $\le$ t) 양측' 값을 확인합니다. 이 값은 약 0.0394입니다.
    \item \textbf{유의 수준 ($\alpha$):} 0.05
    \item \textbf{결론:} P-값 (0.0394)이 유의 수준 $\alpha$ (0.05)보다 작습니다. ($0.0394 < 0.05$)
\end{itemize}
따라서, 우리는 귀무 가설을 기각합니다. 이는 5\% 유의 수준에서 두 제품의 평균 평가 점수가 통계적으로 유의미하게 다르다는 것을 의미합니다.

\subsection{어떤 제품이 더 우수한가?}
Excel 출력에서 제품 A의 평균 점수(61.752)가 제품 B의 평균 점수(57.312)보다 높습니다. 귀무 가설을 기각했으므로, 이 차이는 우연이 아니며 제품 A가 제품 B보다 더 높은 평가를 받았다고 결론 내릴 수 있습니다.

\begin{examplebox}[title=단측 검정으로 가설을 설정했다면?]
만약 처음부터 제품 A가 더 우수할 것이라고 믿고 다음과 같이 가설을 설정했다면:
\begin{itemize}
    \item $H_0: \mu_d \le 0$ (제품 A와 B의 차이가 0이거나 작다)
    \item $H_A: \mu_d > 0$ (제품 A와 B의 차이가 0보다 크다, 즉 A가 B보다 우수하다)
\end{itemize}
이 경우, 'P(T $\le$ t) 단측' 값인 0.0197을 사용합니다. 이 값은 $\alpha=0.05$보다 작으므로 귀무 가설을 기각하고, 제품 A가 제품 B보다 우수하다는 결론을 내릴 수 있습니다.
\end{examplebox}

\subsection{얼마나 더 우수한가? (신뢰 구간 계산)}
두 제품 평가 점수의 평균 차이에 대한 신뢰 구간을 계산하여 제품 A가 얼마나 더 우수한지 정량화할 수 있습니다. 짝지은 표본 검정의 신뢰 구간 공식은 다음과 같습니다.
$\bar{d} \pm t_{\alpha/2} \frac{S_d}{\sqrt{n}}$
여기서,
\begin{itemize}
    \item $\bar{d}$: 쌍별 차이의 평균
    \item $t_{\alpha/2}$: 자유도 $n-1$과 유의 수준 $\alpha$에 해당하는 t-값
    \item $S_d$: 쌍별 차이의 표준 편차
    \item $n$: 쌍의 수 (관측치 수)
\end{itemize}

\begin{enumerate}
    \item \textbf{평균 차이 ($\bar{d}$):}
    Excel 출력에서 제품 A 평균 - 제품 B 평균 = $61.752 - 57.312 = 4.44$
    \item \textbf{t-값 ($t_{\alpha/2}$):}
    90\% 신뢰 구간을 계산하기 위해 $\alpha = 0.10$ (양측)을 사용합니다. 이때의 t-값은 'T 임계값 단측' (0.05 유의 수준)인 1.65를 사용합니다. (단측 0.05는 양측 0.10에 해당합니다.)
    \item \textbf{쌍별 차이의 표준 편차 ($S_d$):}
    이 값은 Excel 출력에서 직접 제공되지 않으므로, 각 참가자의 제품 A 점수와 제품 B 점수의 차이를 계산한 후, 이 차이들의 표준 편차를 별도로 계산해야 합니다. (강의에서는 33.901로 계산되었다고 언급됩니다.)
    \item \textbf{표준 오차:}
    표준 오차 = $\frac{S_d}{\sqrt{n}} = \frac{33.901}{\sqrt{250}} = \frac{33.901}{15.811} \approx 2.144$
    \item \textbf{오차 한계 (Margin of Error):}
    오차 한계 = t-값 $\times$ 표준 오차 = $1.65 \times 2.144 \approx 3.54$
    \item \textbf{90\% 신뢰 구간:}
    하한 = $\bar{d} - \text{오차 한계} = 4.44 - 3.54 = 0.90$
    상한 = $\bar{d} + \text{오차 한계} = 4.44 + 3.54 = 7.98$
\end{enumerate}

\begin{summarybox}[title=90\% 신뢰 구간 해석]
우리는 두 제품의 평가 점수 차이가 0.90점에서 7.98점 사이에 있을 것이라고 90\% 신뢰합니다. 즉, 제품 A가 제품 B보다 평균적으로 최소 0.90점에서 최대 7.98점 더 높은 평가를 받을 것으로 예상됩니다.
\end{summarybox}

\subsection{연습 문제: 교육 프로그램 효과 검정}
\begin{examplebox}[title=시나리오:]
학생들이 준비 수업을 이수한 후 표준화된 시험에서 더 나은 성과를 낼 수 있는지 알아보기 위해 65명의 학생 그룹에게 교육 전후로 시험을 치르게 했습니다. 5\% 유의 수준에서 이 주장을 검정합니다. 귀무 가설과 대립 가설은 무엇일까요?
\end{examplebox}

\begin{itemize}
    \item \textbf{짝지은 표본:} 동일한 학생들을 대상으로 교육 전후를 비교하므로 짝지은 표본 검정입니다.
    \item \textbf{주장:} 학생들이 교육 후 '더 나은' 성과를 낼 것이다. 즉, 점수가 증가할 것이다.
    \item \textbf{가설 설정:}
    \begin{itemize}
        \item $H_0: \mu_d \le 0$ (교육 후 점수 - 교육 전 점수 $\le 0$, 즉 점수 향상이 없거나 감소)
        \item $H_A: \mu_d > 0$ (교육 후 점수 - 교육 전 점수 $> 0$, 즉 점수 향상이 있다)
    \end{itemize}
    \item \textbf{유의 수준 ($\alpha$):} 0.05
\end{itemize}

\subsection{연습 문제: Excel 결과 및 결론}
Excel의 't-검정: 짝지은 두 표본 평균' 분석 결과는 다음과 같습니다.

\begin{table}[h!]
\caption{교육 전/후 t-검정: 짝지은 두 표본 평균}
\label{tab:training_ttest_paired}
\centering
\begin{adjustbox}{width=\textwidth,center}
\begin{tabular}{lcc}
\toprule
\textbf{항목} & \textbf{교육 전 점수} & \textbf{교육 후 점수} \\
\midrule
평균 (Mean) & 500.9692308 & 548.6 \\
분산 (Variance) & 7243.936538 & 6584.93125 \\
관측치 (Observations) & 65 & 65 \\
Pearson 상관 계수 (Pearson Correlation) & 0.099719694 & \\
가설 평균 차이 (Hypothesized Mean Difference) & 0 & \\
자유도 (df) & 64 & \\
t-통계량 (t-stat) & -3.441396891 & \\
P(T $\le$ t) 단측 (P(T $\le$ t) one-tail) & \textbf{0.000512213} & \\
T 임계값 단측 (T critical one-tail) & 1.669013025 & \\
P(T $\le$ t) 양측 (P(T $\le$ t) two-tail) & 0.001024426 & \\
T 임계값 양측 (T critical two-tail) & 1.997729654 & \\
\bottomrule
\end{tabular}
\end{adjustbox}
\end{table}

\begin{itemize}
    \item \textbf{P-값 (단측):} 우리의 가설은 $H_A: \mu_d > 0$ (교육 후 점수 - 교육 전 점수)이므로, Excel 출력에서 'P(T $\le$ t) 단측' 값을 확인합니다. (여기서 t-stat이 음수이므로, 실제로는 $P(T \ge |t|)$ 또는 $P(T \le t)$의 반대 꼬리를 봐야 합니다. Excel은 보통 t-stat의 부호에 따라 적절한 꼬리 확률을 제공합니다. 강의에서는 0.0005를 사용하므로 이를 따릅니다.) 이 값은 0.000512213입니다.
    \item \textbf{유의 수준 ($\alpha$):} 0.05
    \item \textbf{결론:} P-값 (0.000512213)이 유의 수준 $\alpha$ (0.05)보다 작습니다. ($0.000512213 < 0.05$)
\end{itemize}
따라서, 귀무 가설을 기각합니다. 이 데이터에 따르면, 교육 프로그램이 학생들의 시험 점수를 통계적으로 유의미하게 향상시킨다고 결론 내릴 수 있습니다.

\begin{summarybox}[title=짝지은 표본 검정의 장점]
\begin{itemize}
    \item \textbf{개체 간 변동성 제거:} 동일한 개체를 사용하거나 특성을 짝지음으로써, 개체 자체의 차이로 인한 변동성을 제거하고 두 처리 또는 시점 간의 순수한 차이에 집중할 수 있습니다.
    \item \textbf{더 작은 표본 크기:} 변동성이 줄어들기 때문에, 독립 표본 검정보다 더 적은 수의 표본으로도 통계적으로 유의미한 결과를 얻을 수 있습니다.
    \item \textbf{더 신뢰할 수 있는 결과:} 변동성 제어 덕분에 분석 결과의 신뢰도가 높아집니다.
\end{itemize}
가능하다면 항상 짝지은 표본 검정을 선택하는 것이 좋습니다.
\end{summarybox}

\newpage
\section{Lesson 2-4.2: Excel을 이용한 짝지은 표본 검정}

\subsection{Excel을 이용한 짝지은 표본 검정 절차}
이전 연습 문제(교육 프로그램 효과 검정)를 Excel에서 직접 수행하는 방법을 단계별로 살펴보겠습니다.

\begin{examplebox}[title=시나리오:]
학생들이 준비 수업을 이수한 후 표준화된 시험에서 더 나은 성과를 낼 수 있는지 알아보기 위해 65명의 학생 그룹에게 교육 전후로 시험을 치르게 했습니다. 5\% 유의 수준에서 이 주장을 검정합니다.
\end{examplebox}

\subsection{가설 설정}
\begin{itemize}
    \item \textbf{귀무 가설 ($H_0$):} 교육 후 점수와 교육 전 점수의 평균 차이($\mu_d$)는 0보다 작거나 같다. ($\mu_d \le 0$)
    \item \textbf{대립 가설 ($H_A$):} 교육 후 점수와 교육 전 점수의 평균 차이($\mu_d$)는 0보다 크다. ($\mu_d > 0$)
    \item \textbf{유의 수준 ($\alpha$):} 0.05
\end{itemize}

\subsection{Excel 분석 단계}
\begin{enumerate}
    \item \textbf{데이터 준비:} 교육 전 점수(Pre-training score)와 교육 후 점수(Post-training score)를 각각 다른 열에 입력합니다. (예: B열에 교육 전, C열에 교육 후)
    \item \textbf{데이터 분석 도구 실행:}
    \begin{itemize}
        \item '데이터' 탭 $\rightarrow$ '데이터 분석'을 선택합니다.
        \item 't-검정: 짝지은 두 표본 평균 (t-Test: Paired Two Sample for Means)'을 선택하고 '확인'을 클릭합니다.
    \end{itemize}
    \item \textbf{입력 설정:}
    \begin{itemize}
        \item \textbf{변수 1 범위 (Variable 1 Range):} 교육 전 점수 열을 선택합니다 (예: \$B\$B).
        \item \textbf{변수 2 범위 (Variable 2 Range):} 교육 후 점수 열을 선택합니다 (예: \$C\$C).
        \item \textbf{가설 평균 차이 (Hypothesized Mean Difference):} 0을 입력합니다. (우리는 차이가 0보다 큰지 검정하므로, 귀무 가설의 기준점은 0입니다.)
        \item \textbf{레이블 (Labels):} 열 제목이 포함되어 있다면 선택합니다.
        \item \textbf{알파 (Alpha):} 0.05를 입력합니다.
        \item \textbf{출력 옵션 (Output Options):} 결과를 표시할 위치를 선택합니다 (예: '출력 범위'를 선택하고 특정 셀을 지정).
    \end{itemize}
    \item \textbf{결과 확인:} '확인'을 클릭하여 분석 결과를 얻습니다.
\end{enumerate}

\subsection{Excel 결과 해석}
Excel 출력 결과는 다음과 같습니다.

\begin{table}[h!]
\caption{교육 전/후 t-검정: 짝지은 두 표본 평균 (Excel 출력)}
\label{tab:training_ttest_excel}
\centering
\begin{adjustbox}{width=\textwidth,center}
\begin{tabular}{lcc}
\toprule
\textbf{항목} & \textbf{교육 전 점수} & \textbf{교육 후 점수} \\
\midrule
평균 (Mean) & 500.9692308 & 548.6 \\
분산 (Variance) & 7243.936538 & 6584.93125 \\
관측치 (Observations) & 65 & 65 \\
Pearson 상관 계수 (Pearson Correlation) & 0.099719694 & \\
가설 평균 차이 (Hypothesized Mean Difference) & 0 & \\
자유도 (df) & 64 & \\
t-통계량 (t-stat) & -3.441396891 & \\
P(T $\le$ t) 단측 (P(T $\le$ t) one-tail) & \textbf{0.000512213} & \\
T 임계값 단측 (T critical one-tail) & 1.669013025 & \\
P(T $\le$ t) 양측 (P(T $\le$ t) two-tail) & 0.001024426 & \\
T 임계값 양측 (T critical two-tail) & 1.997729654 & \\
\bottomrule
\end{tabular}
\end{adjustbox}
\end{table}

\begin{itemize}
    \item \textbf{P-값 (단측):} 우리의 대립 가설은 $\mu_d > 0$ (교육 후 - 교육 전)이므로, 'P(T $\le$ t) 단측' 값을 확인합니다. (t-stat이 음수이므로, Excel은 보통 t-stat의 절대값에 대한 오른쪽 꼬리 확률을 제공하거나, t-stat의 부호에 따라 적절한 꼬리 확률을 제공합니다. 여기서는 0.000512213을 사용합니다.)
    \item \textbf{유의 수준 ($\alpha$):} 0.05
    \item \textbf{결론:} P-값 (0.000512213)이 유의 수준 $\alpha$ (0.05)보다 훨씬 작습니다. ($0.000512213 < 0.05$)
\end{itemize}
따라서, 귀무 가설을 기각합니다. 이는 교육이 학생들의 시험 점수를 통계적으로 유의미하게 증가시킨다는 것을 의미합니다.

\subsection{신뢰 구간 계산 (평균 차이)}
평균 차이에 대한 신뢰 구간을 계산하여 교육 효과의 크기를 추정할 수 있습니다.
신뢰 구간 공식: $\bar{d} \pm t_{\alpha/2} \frac{S_d}{\sqrt{n}}$

\begin{enumerate}
    \item \textbf{평균 차이 ($\bar{d}$):}
    교육 후 평균 - 교육 전 평균 = $548.6 - 500.9692308 = 47.63076923$
    \item \textbf{t-값 ($t_{\alpha/2}$):}
    단측 검정에서 $\alpha = 0.05$를 사용했으므로, 양측 신뢰 구간을 계산할 때는 양쪽 꼬리에 각각 0.05씩, 즉 총 0.10을 제외한 90\% 신뢰 구간을 얻게 됩니다. 이때 사용되는 t-값은 'T 임계값 단측'인 1.669013025입니다.
    \item \textbf{쌍별 차이의 표준 편차 ($S_d$):}
    이 값은 Excel 출력에서 직접 제공되지 않습니다. 다음 단계에 따라 수동으로 계산해야 합니다.
    \begin{itemize}
        \item \textbf{차이 열 생성:} Excel 데이터 옆에 'diff'라는 새 열을 만듭니다. 각 학생에 대해 '교육 후 점수 - 교육 전 점수'를 계산하여 이 열을 채웁니다. (예: \texttt{=C2-B2})
        \item \textbf{표준 편차 계산:} 'diff' 열의 모든 값에 대해 \texttt{STDEV.S()} 함수를 사용하여 표준 편차를 계산합니다. (예: \texttt{=STDEV.S(D2:D66)})
        \item 계산 결과: $S_d \approx 111.5859491$
    \end{itemize}
    \item \textbf{표준 오차:}
    표준 오차 = $\frac{S_d}{\sqrt{n}} = \frac{111.5859491}{\sqrt{65}} = \frac{111.5859491}{8.062257748} \approx 13.8400404$
    \item \textbf{오차 한계 (Margin of Error):}
    오차 한계 = t-값 $\times$ 표준 오차 = $1.669013025 \times 13.8400404 \approx 23.10003082$
    \item \textbf{90\% 신뢰 구간:}
    하한 (Lower Bound) = 평균 차이 - 오차 한계 = $47.63076923 - 23.10003082 = 24.53073841$
    상한 (Upper Bound) = 평균 차이 + 오차 한계 = $47.63076923 + 23.10003082 = 70.73080005$
\end{enumerate}

\begin{table}[h!]
\caption{통계적 유의성 - 교육 프로그램 (90\% 신뢰 구간)}
\label{tab:training_ci}
\centering
\begin{adjustbox}{width=0.6\textwidth,center}
\begin{tabular}{lc}
\toprule
\textbf{항목} & \textbf{값} \\
\midrule
평균 차이 (mean diff) & 47.63076923 \\
t-값 (t value) & 1.669013025 \\
차이의 표준 편차 (std dev of diff) & 111.5859491 \\
오차 한계 (Margin of error) & 23.10003082 \\
하한 (Lower) & 24.53073841 \\
상한 (Upper) & 70.73080005 \\
\bottomrule
\end{tabular}
\end{adjustbox}
\end{table}

\begin{summarybox}[title=90\% 신뢰 구간 해석]
우리는 이 교육을 받은 학생 모집단의 평균 점수가 24점에서 70점 사이로 증가할 것이라고 90\% 신뢰합니다. 이는 최소 24점, 최대 70점까지 점수가 향상될 수 있다는 의미입니다. 이 범위 내의 모든 값은 그럴듯한(plausible) 값입니다.
\end{summarybox}

\newpage
\section{Lesson 2-5.1: 비율 검정}

\subsection{개념: 두 모집단 비율 비교}
우리는 종종 두 모집단에서 특정 특성을 가진 개체의 '비율(proportion)'을 비교해야 할 때가 있습니다. 예를 들어, 두 지역의 실업률(실업자 비율), 두 후보에 대한 지지율(지지자 비율), 또는 두 공정의 불량률(불량품 비율) 등을 비교할 수 있습니다. 비율 검정은 이러한 비율들이 통계적으로 유의미하게 다른지 여부를 판단하는 데 사용됩니다.

\begin{examplebox}[title=비율 비교의 실제 사례]
\begin{itemize}
    \item \textbf{실업률:} 2016년 4월 실업률(5\%)이 2015년 4월 실업률(5.4\%)과 유의미하게 다른가?
    \item \textbf{정치 여론조사:} 힐러리 클린턴 지지율 40\%, 버니 샌더스 지지율 46\% (오차 한계 $\pm$ 4\%). 샌더스가 확실히 이길 것으로 예측할 수 있는가?
    \begin{itemize}
        \item 클린턴: $40\% \pm 4\% = [36\%, 44\%]$
        \item 샌더스: $46\% \pm 4\% = [42\%, 50\%]$
    \end{itemize}
    이 경우, 샌더스의 최저 지지율(42\%)이 클린턴의 최고 지지율(44\%)보다 낮을 수 있으므로, 샌더스가 확실히 이긴다고 단정하기 어렵습니다. (강의 내용과 달리, 샌더스 42%가 클린턴 44%보다 낮으므로, 이 예시에서는 샌더스가 이긴다고 단정하기 어렵습니다. 강의에서는 샌더스가 이길 수 있다고 언급되었으나, 이는 오차 한계의 해석에 따라 달라질 수 있습니다.)
\end{itemize}
\end{examplebox}

\subsection{비율 검정의 가설}
두 독립적인 모집단에서 추출된 두 표본의 비율을 비교할 때의 가설은 다음과 같습니다.

\begin{table}[h!]
\caption{비율 검정의 가설 유형}
\label{tab:proportion_hypotheses}
\centering
\begin{adjustbox}{width=0.8\textwidth,center}
\begin{tabular}{lll}
\toprule
\textbf{$H_0$} & \textbf{$H_A$} & \textbf{유형} \\
\midrule
$p_1 = p_2$ & $p_1 \ne p_2$ & 양측 검정 (Two-tail) \\
$p_1 \ge p_2$ & $p_1 < p_2$ & 단측 검정 (One-tail) \\
$p_1 \le p_2$ & $p_1 > p_2$ & 단측 검정 (One-tail) \\
\bottomrule
\end{tabular}
\end{adjustbox}
\end{table}
여기서 $p_1$과 $p_2$는 두 모집단의 실제 비율을 나타냅니다.

\subsection{예시 1: 남학생과 여학생의 수학 성적 비교}
\begin{examplebox}[title=시나리오:]
한 연구에 따르면 남학생과 여학생은 수학 시험에서 동등하게 잘 수행한다고 합니다. 한 학군에서 11학년 남학생 400명과 여학생 360명을 무작위로 표본 추출하여 최근 표준화된 시험 점수를 확인했습니다. 남학생 315명과 여학생 280명이 능숙도 수준에 도달했습니다. 5\% 유의 수준에서 학군은 남학생과 여학생이 수학 시험에서 동등하게 잘 수행한다고 주장할 수 있을까요?
\end{examplebox}

\begin{itemize}
    \item \textbf{표본 1 (남학생):} $n_1 = 400$, 능숙도 학생 수 = 315 $\implies \hat{p}_1 = 315/400 = 0.7875$ (78.75\%)
    \item \textbf{표본 2 (여학생):} $n_2 = 360$, 능숙도 학생 수 = 280 $\implies \hat{p}_2 = 280/360 \approx 0.7778$ (77.78\%)
    \item \textbf{가설 설정:} 동등하게 잘 수행하는지 여부를 검정하므로 양측 검정입니다.
    \begin{itemize}
        \item $H_0: p_1 = p_2$ (남학생과 여학생의 능숙도 비율은 같다)
        \item $H_A: p_1 \ne p_2$ (남학생과 여학생의 능숙도 비율은 다르다)
    \end{itemize}
    \item \textbf{유의 수준 ($\alpha$):} 0.05
\end{itemize}

\subsection{P-값 계산 (Z-검정)}
비율 검정에서는 Z-통계량을 사용합니다.
\begin{itemize}
    \item \textbf{Z-통계량 공식:}
    $z = \frac{(\hat{p}_1 - \hat{p}_2) - D_0}{\sigma_{\hat{p}_1 - \hat{p}_2}}$
    여기서 $D_0$는 가설 평균 차이 (여기서는 0)입니다.
    \item \textbf{표준 오차 공식:}
    $\sigma_{\hat{p}_1 - \hat{p}_2} = \sqrt{\frac{p_1(1-p_1)}{n_1} + \frac{p_2(1-p_2)}{n_2}}$
    (귀무 가설 하에서는 $p_1=p_2=p_{pooled}$로 가정하고 통합 비율을 사용합니다.)
\end{itemize}
Excel에는 비율 검정을 위한 내장 함수가 없지만, 사용자 정의 워크시트를 통해 결과를 얻을 수 있습니다.

\subsection{Excel 출력 및 결론 (예시 1)}
사용자 정의 Excel 워크시트의 출력은 다음과 같습니다.

\begin{table}[h!]
\caption{남학생/여학생 비율 검정 Excel 출력}
\label{tab:math_proportion_excel}
\centering
\begin{adjustbox}{width=\textwidth,center}
\begin{tabular}{lc}
\toprule
\textbf{항목} & \textbf{값} \\
\midrule
표본 1 비율 (Sample 1 Proportion) & 78.75\% \\
표본 2 비율 (Sample 2 Proportion) & 77.78\% \\
비율 차이 (Proportion Difference) & 0.97\% \\
Z (단측) (Z (One-Tail)) & 1.64485 \\
Z (양측) (Z (Two-Tail)) & 1.95996 \\
표준 오차 (Standard Error) & 0.02997 \\
가설 차이 (Hypothesized Difference) & 0 \\
검정 통계량 (Z-Test) 단측 ($H_0: p_1 - p_2 \ge 0$) & 0.324350049 \\
P-값 (단측) & 0.62716 \\
검정 통계량 (Z-Test) 단측 ($H_0: p_1 - p_2 \le 0$) & 0.324350049 \\
P-값 (단측) & 0.37284 \\
검정 통계량 (Z-Test) 양측 ($H_0: p_1 - p_2 = 0$) & 0.32460345 \\
P-값 (양측) & \textbf{0.74548} \\
\bottomrule
\end{tabular}
\end{adjustbox}
\end{table}

\begin{itemize}
    \item \textbf{P-값 (양측):} 우리의 가설은 양측 검정이므로, 'P-값 (양측)'인 0.74548을 확인합니다.
    \item \textbf{유의 수준 ($\alpha$):} 0.05
    \item \textbf{결론:} P-값 (0.74548)이 유의 수준 $\alpha$ (0.05)보다 큽니다. ($0.74548 > 0.05$)
\end{itemize}
따라서, 귀무 가설을 기각하지 않습니다. 이는 이 학군의 남학생과 여학생 표본이 수학 시험에서 동등하게 잘 수행한다는 것을 의미합니다. 관측된 0.97\%의 차이는 통계적으로 유의미하지 않으며, 우연에 의한 것으로 볼 수 있습니다.

\subsection{예시 2: 패스트푸드 워크플로우 재설계}
\begin{examplebox}[title=시나리오:]
한 패스트푸드 체인이 워크플로우 재설계를 테스트 중입니다. 재설계가 효과적이라면 모든 지점에 적용될 것입니다. 새로운 디자인을 표준으로 채택하려면, 고객 주문이 4분 이내에 완료되는 비율이 현재보다 20\% 더 높아야 합니다.
\end{examplebox}

\begin{itemize}
    \item \textbf{데이터:}
    \begin{itemize}
        \item \textbf{현재 디자인:} 10,000건 중 7,600건이 4분 이내 완료 $\implies \hat{p}_2 = 7600/10000 = 0.76$ (76\%)
        \item \textbf{새로운 디자인:} 6,000건 중 5,400건이 4분 이내 완료 $\implies \hat{p}_1 = 5400/6000 = 0.90$ (90\%)
    \end{itemize}
    \item \textbf{가설 설정:} 새로운 디자인이 현재보다 20\% '더 높아야' 하므로, 이는 단측 검정이며, 가설 평균 차이가 0.20입니다.
    \begin{itemize}
        \item $p_1$: 새로운 디자인의 4분 이내 완료 비율
        \item $p_2$: 현재 디자인의 4분 이내 완료 비율
        \item $H_0: p_1 - p_2 \le 0.20$ (새로운 디자인의 비율 증가가 20\% 이하이다)
        \item $H_A: p_1 - p_2 > 0.20$ (새로운 디자인의 비율 증가가 20\% 초과이다)
    \end{itemize}
    \item \textbf{유의 수준 ($\alpha$):} 0.05
\end{itemize}

\subsection{Excel 출력 및 결론 (예시 2)}
사용자 정의 Excel 워크시트의 출력은 다음과 같습니다. '가설 차이'에 0.20을 입력합니다.

\begin{table}[h!]
\caption{패스트푸드 워크플로우 비율 검정 Excel 출력}
\label{tab:fastfood_proportion_excel}
\centering
\begin{adjustbox}{width=\textwidth,center}
\begin{tabular}{lc}
\toprule
\textbf{항목} & \textbf{값} \\
\midrule
표본 1 비율 (Sample 1 Proportion) & 90\% \\
표본 2 비율 (Sample 2 Proportion) & 76\% \\
비율 차이 (Proportion Difference) & 14\% \\
Z (단측) (Z (One-Tail)) & 1.6449 \\
Z (양측) (Z (Two-Tail)) & 1.96 \\
표준 오차 (Standard Error) & 0.00577 \\
가설 차이 (Hypothesized Difference) & \textbf{0.2000} \\
검정 통계량 (Z-Test) 단측 ($H_0: p_1 - p_2 \ge 0$) & -10.4069 \\
P-값 (단측) & 0 \\
검정 통계량 (Z-Test) 단측 ($H_0: p_1 - p_2 \le 0$) & -10.4069 \\
P-값 (단측) & \textbf{1} \\
검정 통계량 (Z-Test) 양측 ($H_0: p_1 - p_2 = 0$) & 9.4136 \\
P-값 (양측) & 0 \\
\bottomrule
\end{tabular}
\end{adjustbox}
\end{table}

\begin{itemize}
    \item \textbf{P-값 (단측):} 우리의 대립 가설은 $H_A: p_1 - p_2 > 0.20$입니다. Excel 출력에서 $H_0: p_1 - p_2 \le 0$에 해당하는 P-값은 1입니다. (Z-통계량이 음수이므로, $p_1 - p_2$가 0.20보다 작다는 방향의 확률이 매우 높다는 의미입니다.)
    \item \textbf{유의 수준 ($\alpha$):} 0.05
    \item \textbf{결론:} P-값 (1)이 유의 수준 $\alpha$ (0.05)보다 큽니다. ($1 > 0.05$)
\end{itemize}
따라서, 귀무 가설을 기각하지 않습니다. 이는 새로운 디자인이 고객 주문 완료 비율을 최소 20\% 이상 개선한다는 것을 보여주지 못했음을 의미합니다.

\subsection{95\% 신뢰 구간 계산 (비율 차이)}
경영진은 새로운 디자인의 비율 증가에 대한 95\% 신뢰 구간을 알고 싶어 합니다.
신뢰 구간 공식: $(\hat{p}_1 - \hat{p}_2) \pm z_{\alpha/2} \sqrt{\frac{\hat{p}_1(1-\hat{p}_1)}{n_1} + \frac{\hat{p}_2(1-\hat{p}_2)}{n_2}}$

\begin{enumerate}
    \item \textbf{표본 비율 차이 ($\hat{p}_1 - \hat{p}_2$):}
    Excel 출력에서 14\% (0.14)
    \item \textbf{Z-값 ($z_{\alpha/2}$):}
    95\% 신뢰 구간을 위한 Z-값은 'Z (양측)'인 1.96입니다.
    \item \textbf{표준 오차:}
    Excel 출력에서 0.00577
    \item \textbf{신뢰 구간 계산:}
    $[0.14 \pm 1.96 \times 0.00577]$
    $[0.14 \pm 0.0113]$
    $[0.1287, 0.1513]$ 또는 $[12.87\%, 15.13\%]$
\end{enumerate}

\begin{summarybox}[title=95\% 신뢰 구간 해석]
우리는 새로운 디자인이 고객 주문 완료 비율을 현재 디자인보다 12.87\%에서 15.13\% 사이로 증가시킬 것이라고 95\% 신뢰합니다. 경영진은 20\% 이상의 개선을 원했지만, 우리의 연구 결과에 따르면 최대 개선 폭은 약 15\%에 불과합니다.
\end{summarybox}

\subsection{표본 크기의 중요성}
비율 비교에서도 표본 크기는 매우 중요합니다. 표본 크기가 작으면 오차 한계가 커져 통계적으로 유의미한 차이를 발견하기 어렵습니다. 반대로 표본 크기가 크면 작은 차이도 통계적으로 유의미하게 나타날 수 있습니다.

\begin{examplebox}[title=정치 여론조사 사례:]
\begin{itemize}
    \item \textbf{초기 결과 (108명 투표):} 힐러리 클린턴 56.5\%, 버니 샌더스 41.7\%. 클린턴이 15\% 앞서지만, 표본이 너무 작아 오차 한계가 매우 높으므로 승리를 단정할 수 없습니다.
    \item \textbf{최종 결과 (185,000명 투표):} 버니 샌더스 61.7\%, 힐러리 클린턴 38.3\%. 표본 크기가 매우 커지자 샌더스와 클린턴 간의 차이가 통계적으로 유의미해졌고, 샌더스의 승리를 확실히 예측할 수 있습니다.
\end{itemize}
겉보기에 큰 차이도 표본 크기가 작으면 통계적으로 무의미할 수 있습니다. 비율을 비교할 때는 항상 표본 크기를 염두에 두어야 합니다.
\end{examplebox}

\subsection{연습 문제: UI 입학 허가 수락률 감소 주장 검정}
\begin{examplebox}[title=시나리오:]
지역 신문에 "전반적으로, 올해 봄 UI의 입학 허가 수락 학생 수는 작년보다 116명 감소했습니다. 2015년 8,085명에서 2016년 7,969명으로 1\% 감소했습니다."라는 기사가 실렸습니다. 기자는 이 감소가 통계적으로 유의미한지 주장합니다.
\end{examplebox}

\begin{table}[h!]
\caption{UI 입학 허가 수락 데이터}
\label{tab:ui_admission_data}
\centering
\begin{adjustbox}{width=0.6\textwidth,center}
\begin{tabular}{lcc}
\toprule
\textbf{범주} & \textbf{2015년} & \textbf{2016년} \\
\midrule
지원자 수 (Number Applied) & 22,639 & 22,711 \\
UI 제안 수락 (Accepted UI's offer) & 8,085 & 7,969 \\
\bottomrule
\end{tabular}
\end{adjustbox}
\end{table}

\begin{itemize}
    \item \textbf{비율 계산:}
    \begin{itemize}
        \item $p_1$: 2015년 입학 허가 수락 학생 비율 = $8085 / 22639 \approx 0.3571$ (35.71\%)
        \item $p_2$: 2016년 입학 허가 수락 학생 비율 = $7969 / 22711 \approx 0.3509$ (35.09\%)
    \end{itemize}
    \item \textbf{기자의 주장:} 2015년 비율이 2016년 비율보다 '높다' (감소했다).
    \item \textbf{가설 설정:}
    \begin{itemize}
        \item $H_0: p_1 \le p_2$ (2015년 비율이 2016년 비율보다 작거나 같다)
        \item $H_A: p_1 > p_2$ (2015년 비율이 2016년 비율보다 크다)
    \end{itemize}
    \item \textbf{유의 수준 ($\alpha$):} 0.05 (가정)
\end{itemize}

\subsection{연습 문제: Excel 출력 및 결론}
사용자 정의 Excel 워크시트의 출력은 다음과 같습니다.

\begin{table}[h!]
\caption{UI 입학 허가 수락률 비율 검정 Excel 출력}
\label{tab:ui_admission_proportion_excel}
\centering
\begin{adjustbox}{width=\textwidth,center}
\begin{tabular}{lc}
\toprule
\textbf{항목} & \textbf{값} \\
\midrule
표본 1 비율 (Sample 1 Proportion) & 35.71\% \\
표본 2 비율 (Sample 2 Proportion) & 35.09\% \\
비율 차이 (Proportion Difference) & 0.62\% \\
Z (단측) (Z (One-Tail)) & 1.64485 \\
Z (양측) (Z (Two-Tail)) & 1.95996 \\
표준 오차 (Standard Error) & 0.00449 \\
가설 차이 (Hypothesized Difference) & 0 \\
검정 통계량 (Z-Test) 단측 ($H_0: p_1 - p_2 \ge 0$) & 1.389376636 \\
P-값 (단측) & 0.91764 \\
검정 통계량 (Z-Test) 단측 ($H_0: p_1 - p_2 \le 0$) & 1.389376636 \\
P-값 (단측) & \textbf{0.08236} \\
검정 통계량 (Z-Test) 양측 ($H_0: p_1 - p_2 = 0$) & 1.389355854 \\
P-값 (양측) & 0.16472 \\
\bottomrule
\end{tabular}
\end{adjustbox}
\end{table}

\begin{itemize}
    \item \textbf{P-값 (단측):} 우리의 대립 가설은 $H_A: p_1 > p_2$입니다. Excel 출력에서 $H_0: p_1 - p_2 \le 0$에 해당하는 P-값은 0.08236입니다.
    \item \textbf{유의 수준 ($\alpha$):} 0.05
    \item \textbf{결론:} P-값 (0.08236)이 유의 수준 $\alpha$ (0.05)보다 큽니다. ($0.08236 > 0.05$)
\end{itemize}
따라서, 귀무 가설을 기각하지 않습니다. 이는 2016년 입학 허가 수락률이 2015년과 통계적으로 유의미한 차이를 보이지 않는다는 것을 의미합니다. 관측된 0.62\%의 감소는 통계적 '노이즈' 범위 내에 있으며, 중요하다고 볼 만한 변화가 아닙니다.

\begin{warningbox}[title=통계적 유의성과 뉴스 가치]
이 예시는 우리가 읽거나 듣는 모든 정보가 통계적으로 유의미한 결론을 나타내는 것은 아님을 보여줍니다. 통계적 가설 검정은 무엇이 '유의미한 변화'이고 무엇이 '노이즈'인지를 판단하는 데 도움을 줍니다.
\end{warningbox}

\newpage

\newpage
\section{개요: 두 표본 비율 검정}

이 문서는 '추론 및 예측 통계 모듈 2'의 마지막 부분으로, 두 표본의 비율을 비교하는 통계적 가설 검정 방법에 대해 다룹니다. 특히, 단측 검정(One-Tail Test)과 양측 검정(Two-Tail Test)을 엑셀(Excel)을 활용하여 수행하는 절차와 결과를 해석하는 방법을 상세히 설명합니다.

\begin{itemize}
    \item \textbf{주요 목표:} 두 집단의 특정 사건 발생 비율(예: 고객 주문 완료율, 시험 합격률)에 유의미한 차이가 있는지 통계적으로 판단합니다.
    \item \textbf{핵심 개념:} 귀무가설(Null Hypothesis)과 대립가설(Alternative Hypothesis) 설정, 유의수준(Significance Level), p-값(p-value) 해석, 신뢰구간(Confidence Interval) 계산 및 활용.
    \item \textbf{실습 도구:} 엑셀 워크시트를 사용하여 복잡한 계산 없이 쉽게 검정을 수행합니다.
    \item \textbf{학습 목표:} 주어진 시나리오에서 적절한 가설을 설정하고, 엑셀을 이용해 검정을 수행하며, 그 결과를 바탕으로 합리적인 의사결정을 내릴 수 있도록 합니다.
\end{itemize}

\newpage
\section{용어 정리}

통계적 가설 검정에서 사용되는 주요 용어들을 비전공자도 이해하기 쉽게 정리했습니다.

\begin{adjustbox}{width=\textwidth,center}
\begin{tabular}{lllc}
\toprule
\textbf{용어} & \textbf{쉬운 설명} & \textbf{원어} & \textbf{비고} \\
\midrule
\textbf{비율 검정} & 두 집단에서 어떤 특성을 가진 개체의 비율이 서로 같은지 다른지 통계적으로 확인하는 방법. & Test of Proportions & \\
\textbf{단측 검정} & 한 집단의 비율이 다른 집단보다 '크다' 또는 '작다'와 같이 특정 방향의 차이만 검정하는 방법. & One-Tail Test & 대립가설에 부등호($>$, $<$) 포함 \\
\textbf{양측 검정} & 두 집단의 비율이 '다르다'와 같이 차이의 방향을 특정하지 않고 단순히 같지 않음을 검정하는 방법. & Two-Tail Test & 대립가설에 같지 않음($\neq$) 포함 \\
\textbf{귀무가설} & 통계적으로 증명하고 싶은 것과 반대되는, '차이가 없다'거나 '특정 값과 같다'는 기본 가정. & Null Hypothesis ($H_0$) & 보통 기각하기 위해 검정 \\
\textbf{대립가설} & 귀무가설과 반대되는, '차이가 있다'거나 '특정 값보다 크거나 작다'는 주장. & Alternative Hypothesis ($H_A$) & 우리가 증명하고 싶은 주장 \\
\textbf{유의수준} & 귀무가설이 사실인데도 불구하고 잘못 기각할 확률의 최대 허용치. 보통 5\%($\alpha=0.05$)를 사용. & Significance Level ($\alpha$) & 오류를 허용하는 정도 \\
\textbf{p-값} & 귀무가설이 사실이라고 가정했을 때, 현재 관측된 데이터 또는 그보다 더 극단적인 데이터가 나올 확률. & p-value & p-값이 작을수록 귀무가설 기각의 증거가 강함 \\
\textbf{신뢰구간} & 모수(모집단의 실제 값)가 포함될 것으로 예상되는 구간. & Confidence Interval (CI) & 특정 신뢰수준(예: 95\%)에서 모수를 포함할 확률 \\
\textbf{표본 비율} & 표본에서 특정 특성을 가진 개체의 비율. & Sample Proportion ($\hat{p}$) & 모집단 비율($p$)의 추정치 \\
\textbf{가설 차이} & 귀무가설에서 설정한 두 비율 간의 예상 차이 값. & Hypothesized Difference & 보통 0 또는 특정 값(예: 0.20) \\
\textbf{표준 오차} & 표본 통계량(예: 표본 비율 차이)이 모집단 모수와 얼마나 차이 날 수 있는지 나타내는 척도. & Standard Error (SE) & 추정치의 정밀도 \\
\textbf{z-값} & 표준정규분포에서 특정 확률에 해당하는 값. 신뢰구간 계산 등에 사용. & z-value & 신뢰수준에 따라 달라짐 \\
\textbf{마진 오브 에러} & 신뢰구간의 폭을 결정하는 값. 표본 통계량에 더하고 빼서 신뢰구간을 만듦. & Margin of Error & z-값 $\times$ 표준 오차 \\
\bottomrule
\end{tabular}
\end{adjustbox}

\newpage
\section{Lesson 2-5.2: 두 표본 비율에 대한 단측 검정 (엑셀 활용)}

이 섹션에서는 두 표본의 비율 차이가 특정 값보다 큰지(또는 작은지)를 검정하는 단측 검정 방법을 엑셀을 사용하여 학습합니다. 특히, 엑셀 워크시트를 활용하여 계산 과정을 간소화하고 결과를 해석하는 데 중점을 둡니다.

\subsection{핵심 개념 및 원리}

단측 검정은 우리가 특정 방향의 변화(예: '더 좋아졌다', '더 나빠졌다')를 확인하고 싶을 때 사용합니다. 예를 들어, 새로운 마케팅 전략이 고객 전환율을 '높였는지', 새로운 생산 공정이 불량률을 '낮췄는지' 등을 검정할 때 유용합니다.

\begin{summarybox}[title=단측 검정의 핵심:]
두 비율의 차이가 특정 값보다 크거나 작은지 한 방향으로만 검정합니다.
귀무가설($H_0$)과 대립가설($H_A$)은 부등호($>$, $<$, $\le$, $\ge$)를 포함합니다.
\end{summarybox}

\subsubsection{가설 설정}
단측 검정에서는 우리가 증명하고자 하는 주장이 대립가설($H_A$)이 됩니다. 귀무가설($H_0$)은 그 반대 또는 '차이가 없다'는 가정이 됩니다.

\begin{itemize}
    \item \textbf{귀무가설 ($H_0$):} 두 비율의 차이가 특정 값보다 작거나 같다. (예: $P_1 - P_2 \le 0.20$)
    \item \textbf{대립가설 ($H_A$):} 두 비율의 차이가 특정 값보다 크다. (예: $P_1 - P_2 > 0.20$)
\end{itemize}
또는 반대로:
\begin{itemize}
    \item \textbf{귀무가설 ($H_0$):} 두 비율의 차이가 특정 값보다 크거나 같다. (예: $P_1 - P_2 \ge 0.20$)
    \item \textbf{대립가설 ($H_A$):} 두 비율의 차이가 특정 값보다 작다. (예: $P_1 - P_2 < 0.20$)
\end{itemize}

\begin{cautionbox}[title=가설 설정 시 주의사항:]
어떤 가설을 설정하느냐에 따라 p-값 계산 방식이 달라지므로, 문제의 요구사항을 정확히 파악하여 올바른 귀무가설과 대립가설을 설정하는 것이 매우 중요합니다. 특히, 엑셀 워크시트에서 어떤 단측 검정 p-값을 봐야 할지 결정하는 기준이 됩니다.
\end{cautionbox}

\subsection{패스트푸드 체인점 워크플로우 개선 예시}

패스트푸드 체인점에서 새로운 워크플로우를 테스트하여, 고객 주문 완료율이 4분 이내에 20\% 더 높아지는지 확인하고자 합니다. 만약 효과적이라면 모든 지점에 적용할 계획입니다.

\subsubsection{문제 정의 및 데이터 수집}
\begin{itemize}
    \item \textbf{목표:} 새로운 디자인이 현재 디자인보다 4분 이내 주문 완료 비율을 20\% 이상 높이는지 확인.
    \item \textbf{데이터:}
    \begin{itemize}
        \item \textbf{현재 디자인 (Current Design):} 10,000건 중 7,600건이 4분 이내 완료.
        \item \textbf{새로운 디자인 (New Design):} 6,000건 중 5,400건이 4분 이내 완료.
    \end{itemize}
    \item \textbf{유의수준 ($\alpha$):} 5\% (0.05)
\end{itemize}

\subsubsection{가설 설정}
새로운 디자인(P1)이 현재 디자인(P2)보다 20\% 더 높은 비율을 보이는지 확인하므로, 대립가설은 'P1 - P2 > 0.20'이 됩니다.

\begin{itemize}
    \item $P_1$: 새로운 디자인에서의 4분 이내 주문 완료 비율
    \item $P_2$: 현재 디자인에서의 4분 이내 주문 완료 비율
    \item \textbf{귀무가설 ($H_0$):} $P_1 - P_2 \le 0.20$ (새로운 디자인이 20\% 이상 개선되지 않았다)
    \item \textbf{대립가설 ($H_A$):} $P_1 - P_2 > 0.20$ (새로운 디자인이 20\% 이상 개선되었다)
\end{itemize}
우리는 대립가설($H_A$)을 지지할 증거를 찾고 있습니다.

\subsection{엑셀을 이용한 검정 절차}

제공된 엑셀 워크시트('Proportion Fast Food' 파일의 'Two proportions - First worksheet' 참조)를 사용하여 검정을 수행합니다.

\subsubsection{데이터 입력}
엑셀 워크시트의 'Input' 섹션에 다음 값을 입력합니다.

\begin{itemize}
    \item \textbf{Sample 1 (New Design):}
    \begin{itemize}
        \item Count of Events: 5,400
        \item Sample Size: 6,000
    \end{itemize}
    \item \textbf{Sample 2 (Current Design):}
    \begin{itemize}
        \item Count of Events: 7,600
        \item Sample Size: 10,000
    \end{itemize}
    \item \textbf{Level of Significance ($\alpha$):} 0.05
    \item \textbf{Hypothesized Difference:} 0.20 (우리가 20\% 개선을 가정했으므로)
\end{itemize}

\begin{tipbox}
엑셀 워크시트는 초기에는 오류 메시지를 표시할 수 있습니다. 데이터를 모두 입력하면 'Output' 부분이 자동으로 채워집니다.
\end{tipbox}

\subsubsection{결과 해석 (p-값)}
데이터를 입력하면 엑셀 워크시트의 'Output' 섹션이 계산 결과로 채워집니다. 우리는 설정한 가설에 맞는 p-값을 찾아야 합니다.

\begin{itemize}
    \item \textbf{표본 비율 계산:}
    \begin{itemize}
        \item 새로운 디자인($\hat{p}_1$): $5400 / 6000 = 0.90$ (90\%)
        \item 현재 디자인($\hat{p}_2$): $7600 / 10000 = 0.76$ (76\%)
        \item 표본 비율 차이($\hat{p}_1 - \hat{p}_2$): $0.90 - 0.76 = 0.14$ (14\%)
    \end{itemize}
    \item \textbf{p-값 확인:}
        우리의 귀무가설은 $H_0: P_1 - P_2 \le 0.20$ 이므로, 엑셀 워크시트에서 이 형태에 해당하는 단측 검정(One-tail) p-값을 찾습니다.
        \begin{itemize}
            \item 엑셀 출력에서 $H_0: P_1 - P_2 \le 0$ (또는 이와 유사한 형태, 여기서 0은 가설 차이 0.20을 의미)에 해당하는 p-값은 \textbf{1.000}입니다.
        \end{itemize}
\end{itemize}

\begin{summarybox}[title=p-값 해석:]
\begin{itemize}
    \item p-값 (1.000) $>$ 유의수준 ($\alpha=0.05$)
    \item 따라서, 귀무가설($H_0: P_1 - P_2 \le 0.20$)을 기각할 수 없습니다.
\end{itemize}
\textbf{결론:} 새로운 디자인이 4분 이내 주문 완료 비율을 20\% 이상 개선했다는 통계적 증거를 찾지 못했습니다.
\end{summarybox}

\subsection{신뢰구간 계산 및 활용}

비록 20\% 개선 목표를 달성하지 못했지만, 실제 개선 정도가 어느 정도인지 경영진에게 보여줄 필요가 있습니다. 이때 신뢰구간을 활용할 수 있습니다.

\subsubsection{신뢰구간 공식}
두 비율 차이에 대한 신뢰구간 공식은 다음과 같습니다.
$$ (\hat{p}_1 - \hat{p}_2) \pm z \sqrt{\frac{\hat{p}_1(1-\hat{p}_1)}{n_1} + \frac{\hat{p}_2(1-\hat{p}_2)}{n_2}} $$
여기서:
\begin{itemize}
    \item $(\hat{p}_1 - \hat{p}_2)$: 두 표본 비율의 차이
    \item $z$: 신뢰수준에 해당하는 z-값
    \item $\sqrt{\frac{\hat{p}_1(1-\hat{p}_1)}{n_1} + \frac{\hat{p}_2(1-\hat{p}_2)}{n_2}}$: 표준 오차 (Standard Error, SE)
\end{itemize}

\subsubsection{z-값 선택}
신뢰구간을 계산할 때 사용하는 z-값은 신뢰수준에 따라 달라집니다.
\begin{itemize}
    \item \textbf{단측 검정의 z-값 (예: 90\% 신뢰수준):} 단측 검정에서는 5\%의 유의수준이 한쪽 꼬리에 집중됩니다. 따라서 90\% 신뢰구간을 원한다면 z-값은 1.645를 사용합니다. (양쪽 꼬리에 각각 5\%씩 남는 90\% 신뢰구간과 동일)
    \item \textbf{양측 검정의 z-값 (예: 95\% 신뢰수준):} 일반적으로 95\% 신뢰구간을 계산할 때는 z-값 1.96을 사용합니다. 이는 양쪽 꼬리에 각각 2.5\%씩 남는 경우입니다.
\end{itemize}
강의에서는 95\% 신뢰구간을 계산하기 위해 z-값으로 \textbf{1.96}을 사용합니다. 이는 단측 검정을 수행했더라도 신뢰구간 자체는 양측으로 계산하는 것이 일반적이기 때문입니다.

\begin{cautionbox}[title=z-값 선택의 혼동 방지:]
단측 검정을 수행했더라도 신뢰구간을 계산할 때는 보통 양측 신뢰구간의 z-값을 사용합니다. 예를 들어, 95\% 신뢰구간을 원한다면 z-값 1.96을 사용합니다. 만약 90\% 신뢰구간을 원한다면 z-값 1.645를 사용합니다. 이는 신뢰구간이 '모수가 이 구간 안에 있을 확률'을 나타내기 때문에, 양쪽으로 오차를 고려하는 것이 일반적이기 때문입니다.
\end{cautionbox}

\subsubsection{신뢰구간 계산}
엑셀 워크시트의 'Output'에서 다음 값을 확인합니다.
\begin{itemize}
    \item \textbf{표본 비율 차이 ($\hat{p}_1 - \hat{p}_2$):} 0.14 (14\%)
    \item \textbf{표준 오차 (Standard Error, SE):} 0.006
\end{itemize}
이제 95\% 신뢰구간을 계산합니다.
\begin{itemize}
    \item \textbf{z-값:} 1.96 (95\% 신뢰수준)
    \item \textbf{마진 오브 에러 (Margin of Error):} $z \times SE = 1.96 \times 0.006 \approx 0.01176$ (약 1.18\%)
    \item \textbf{하한 (Lower Bound):} $0.14 - 0.01176 \approx 0.12824$ (12.82\%)
    \item \textbf{상한 (Upper Bound):} $0.14 + 0.01176 \approx 0.15176$ (15.18\%)
\end{itemize}
엑셀 워크시트에서는 반올림하여 다음과 같이 표시됩니다.
\begin{itemize}
    \item Margin of error: 0.0113
    \item 95\% Lower confidence level: 12.87\%
    \item 95\% Upper confidence level: 15.13\%
\end{itemize}

\subsection{최종 결론}

\begin{summarybox}[title=종합 결론:]
\begin{itemize}
    \item 통계적 검정 결과, 새로운 디자인이 4분 이내 주문 완료 비율을 20\% 이상 개선했다는 귀무가설을 기각할 수 없었습니다. 즉, 20\% 개선 목표는 달성되지 않았습니다.
    \item 하지만, 95\% 신뢰구간은 새로운 디자인이 현재 디자인보다 4분 이내 주문 완료 비율을 \textbf{12.87\%에서 15.13\%} 더 높일 가능성이 있다고 보여줍니다.
    \item 이는 경영진에게 새로운 디자인이 비록 20\% 목표에는 미치지 못하지만, 여전히 유의미한 개선 효과(약 13\%~15\%)를 가져올 수 있다는 추가 정보를 제공합니다.
\end{itemize}
\end{summarybox}

\newpage
\section{Lesson 2-5.3: 두 표본 비율에 대한 양측 검정 (엑셀 활용)}

이 섹션에서는 두 표본의 비율이 서로 같은지 다른지를 검정하는 양측 검정 방법을 엑셀을 사용하여 학습합니다. 양측 검정은 차이의 방향을 특정하지 않고 단순히 '차이가 있다/없다'를 확인하고 싶을 때 사용합니다.

\subsection{핵심 개념 및 원리}

양측 검정은 두 집단 간에 어떤 차이가 있는지 여부만 궁금할 때 사용합니다. 예를 들어, 두 가지 교육 방법이 학생들의 성적에 '차이를 만드는지', 또는 두 가지 약물의 부작용 발생률이 '다른지' 등을 검정할 때 활용됩니다.

\begin{summarybox}[title=양측 검정의 핵심:]
두 비율이 서로 같은지 다른지만 검정합니다. 차이의 방향은 중요하지 않습니다.
귀무가설($H_0$)은 등호($=$)를, 대립가설($H_A$)은 같지 않음($\neq$)을 포함합니다.
\end{summarybox}

\subsubsection{가설 설정}
양측 검정의 가설은 항상 다음과 같은 형태를 가집니다.

\begin{itemize}
    \item \textbf{귀무가설 ($H_0$):} 두 비율이 같다. (예: $P_1 = P_2$ 또는 $P_1 - P_2 = 0$)
    \item \textbf{대립가설 ($H_A$):} 두 비율이 같지 않다. (예: $P_1 \neq P_2$ 또는 $P_1 - P_2 \neq 0$)
\end{itemize}
우리는 대립가설($H_A$)을 지지할 증거를 찾고 있습니다. 즉, 두 비율이 다르다는 것을 증명하고 싶습니다.

\subsection{남학생과 여학생의 수학 시험 성적 예시}

한 연구에서 남학생과 여학생이 수학 시험에서 동등하게 잘 수행한다고 주장했습니다. 한 학군에서 11학년 남학생 400명과 여학생 360명을 무작위로 표본 추출하여 최근 표준화된 시험 점수를 조사했습니다.

\subsubsection{문제 정의 및 데이터 수집}
\begin{itemize}
    \item \textbf{목표:} 남학생과 여학생이 수학 시험에서 동등하게 잘 수행한다고 주장할 수 있는지 확인.
    \item \textbf{데이터:}
    \begin{itemize}
        \item \textbf{남학생 (Boys):} 400명 중 315명이 능숙(proficient) 수준.
        \item \textbf{여학생 (Girls):} 360명 중 280명이 능숙(proficient) 수준.
    \end{itemize}
    \item \textbf{유의수준 ($\alpha$):} 5\% (0.05)
\end{itemize}

\subsubsection{가설 설정}
남학생과 여학생이 동등하게 잘 수행하는지(즉, 비율이 같은지) 확인하므로, 양측 검정을 수행합니다.

\begin{itemize}
    \item $p_{Girls}$: 여학생의 능숙 수준 비율
    \item $p_{Boys}$: 남학생의 능숙 수준 비율
    \item \textbf{귀무가설 ($H_0$):} $p_{Girls} = p_{Boys}$ (여학생과 남학생의 능숙 수준 비율이 같다)
    \item \textbf{대립가설 ($H_A$):} $p_{Girls} \neq p_{Boys}$ (여학생과 남학생의 능숙 수준 비율이 같지 않다)
\end{itemize}

\subsection{엑셀을 이용한 검정 절차}

제공된 엑셀 워크시트('Proportion Math' 파일의 'Two proportions - First worksheet' 참조)를 사용하여 검정을 수행합니다.

\subsubsection{데이터 입력}
엑셀 워크시트의 'Input' 섹션에 다음 값을 입력합니다.

\begin{itemize}
    \item \textbf{Sample 1 (Boys):}
    \begin{itemize}
        \item Count of Events: 315
        \item Sample Size: 400
    \end{itemize}
    \item \textbf{Sample 2 (Girls):}
    \begin{itemize}
        \item Count of Events: 280
        \item Sample Size: 360
    \end{itemize}
    \item \textbf{Level of Significance ($\alpha$):} 0.05
    \item \textbf{Hypothesized Difference:} 0 (두 비율이 같다고 가정하므로 차이는 0)
\end{itemize}

\begin{tipbox}
엑셀 워크시트의 노란색으로 강조된 부분에 데이터를 입력하면, 'Output' 섹션이 자동으로 계산 결과로 채워집니다.
\end{tipbox}

\subsubsection{결과 해석 (p-값)}
데이터를 입력하면 엑셀 워크시트의 'Output' 섹션이 계산 결과로 채워집니다.

\begin{itemize}
    \item \textbf{표본 비율 계산:}
    \begin{itemize}
        \item 남학생($\hat{p}_{Boys}$): $315 / 400 = 0.7875$ (78.75\%)
        \item 여학생($\hat{p}_{Girls}$): $280 / 360 \approx 0.7778$ (77.78\%)
        \item 표본 비율 차이($\hat{p}_{Boys} - \hat{p}_{Girls}$): $0.7875 - 0.7778 \approx 0.0097$ (0.97\%)
    \end{itemize}
    \item \textbf{p-값 확인:}
        우리의 귀무가설은 $H_0: P_1 - P_2 = 0$ 이므로, 엑셀 워크시트에서 'Two-tail' 검정에 해당하는 p-값을 찾습니다.
        \begin{itemize}
            \item 엑셀 출력에서 'Two-tail ($H_0: P_1 - P_2 = 0$)'에 해당하는 p-값은 \textbf{0.745}입니다.
        \end{itemize}
\end{itemize}

\begin{summarybox}[title=p-값 해석:]
\begin{itemize}
    \item p-값 (0.745) $>$ 유의수준 ($\alpha=0.05$)
    \item 따라서, 귀무가설($H_0: p_{Girls} = p_{Boys}$)을 기각할 수 없습니다.
\end{itemize}
\textbf{결론:} 남학생과 여학생이 수학 시험에서 동등하게 잘 수행하지 않는다는 통계적 증거를 찾지 못했습니다. 즉, 두 집단 간에 유의미한 차이가 있다고 볼 수 없습니다.
\end{summarybox}

\subsection{최종 결론}

\begin{summarybox}[title=종합 결론:]
\begin{itemize}
    \item 통계적 검정 결과, 남학생과 여학생의 수학 시험 능숙도 비율에 유의미한 차이가 있다는 대립가설을 지지할 증거가 부족합니다.
    \item 따라서, 학군은 남학생과 여학생이 수학 시험에서 \textbf{동등하게 잘 수행한다}고 주장할 수 있습니다.
    \item 관찰된 약 0.97\%의 차이는 통계적으로 유의미하다고 볼 수 없으며, 이는 단순한 표본 오차로 인한 것일 가능성이 높습니다.
\end{itemize}
\end{summarybox}

\newpage
\section{빠르게 훑어보기: 두 표본 비율 검정 요약}

\begin{summarybox}[title=\textbf{두 표본 비율 검정 핵심 요약}]

\textbf{1. 목적:} 두 집단의 비율($P_1, P_2$)이 서로 같은지, 또는 특정 값만큼 차이가 나는지 통계적으로 판단.

\textbf{2. 가설 설정:}
\begin{itemize}
    \item \textbf{단측 검정 (One-Tail Test):} 특정 방향의 차이 검정 (예: $P_1 - P_2 > 0.20$ 또는 $P_1 - P_2 < 0.20$)
        \begin{itemize}
            \item $H_0: P_1 - P_2 \le \text{가설 차이}$
            \item $H_A: P_1 - P_2 > \text{가설 차이}$
        \end{itemize}
    \item \textbf{양측 검정 (Two-Tail Test):} 차이의 유무만 검정 (예: $P_1 - P_2 \neq 0$)
        \begin{itemize}
            \item $H_0: P_1 - P_2 = \text{가설 차이}$
            \item $H_A: P_1 - P_2 \neq \text{가설 차이}$
        \end{itemize}
\end{itemize}

\textbf{3. 데이터 준비:}
\begin{itemize}
    \item 각 표본의 성공 횟수 (Count of Events)
    \item 각 표본의 총 관측치 수 (Sample Size)
    \item 유의수준 ($\alpha$, 보통 0.05)
    \item 가설 차이 (Hypothesized Difference, $H_0$에 설정된 값)
\end{itemize}

\textbf{4. 엑셀 활용 절차:}
\begin{enumerate}
    \item 제공된 엑셀 워크시트 열기.
    \item 'Input' 섹션에 각 표본의 데이터, 유의수준, 가설 차이 입력.
    \item 'Output' 섹션에서 설정한 가설에 맞는 p-값 확인.
\end{enumerate}

\textbf{5. 결과 해석 (p-값):}
\begin{itemize}
    \item \textbf{p-값 $<$ $\alpha$:} 귀무가설($H_0$)을 기각하고 대립가설($H_A$)을 채택합니다. (통계적으로 유의미한 차이/효과가 있음)
    \item \textbf{p-값 $\ge$ $\alpha$:} 귀무가설($H_0$)을 기각할 수 없습니다. (통계적으로 유의미한 차이/효과를 증명하지 못함)
\end{itemize}

\textbf{6. 신뢰구간 (Confidence Interval) 활용 (선택 사항):}
\begin{itemize}
    \item 실제 비율 차이가 어느 범위에 있을지 추정하여 추가적인 정보를 제공합니다.
    \item 공식: $(\hat{p}_1 - \hat{p}_2) \pm z \times \text{SE}$
    \item z-값은 신뢰수준에 따라 선택 (예: 95\% CI는 $z=1.96$)
\end{itemize}

\textbf{7. 최종 의사결정:}
\begin{itemize}
    \item p-값과 신뢰구간을 종합적으로 고려하여 문제에 대한 결론을 내립니다.
    \item 통계적 유의미성뿐만 아니라 실질적인 중요성도 함께 고려해야 합니다.
\end{itemize}
\end{summarybox}

\end{document}
